\documentclass[12pt]{article}

%%%%%%%%%%%%%%%%%%%%%%%%%%%%%%%%%%%%%%%%
% PACKAGES
%%%%%%%%%%%%%%%%%%%%%%%%%%%%%%%%%%%%%%%%
\usepackage[margin=1in]{geometry}
\usepackage[T1]{fontenc}
\usepackage{lmodern}
\usepackage{amsmath,amssymb}
\usepackage{xcolor}
\usepackage{graphicx}
\usepackage{hyperref}
\usepackage{enumitem}
\usepackage{titlesec}
\usepackage{array}
\usepackage{tabularx}
\usepackage{colortbl}
\usepackage{tcolorbox}
\tcbuselibrary{breakable}
\usepackage{listings}
\usepackage[normalem]{ulem}
\lstdefinelanguage{yaml}{
  basicstyle=\ttfamily\small,
  morecomment=[l]{\#},
  morestring=[b]",
  sensitive=true
}

%%%%%%%%%%%%%%%%%%%%%%%%%%%%%%%%%%%%%%%%
% COLOR SCHEME (high-contrast, consistent theme)
%%%%%%%%%%%%%%%%%%%%%%%%%%%%%%%%%%%%%%%%
\definecolor{Navy}{RGB}{16,52,96}
\definecolor{Teal}{RGB}{0,132,105}
\definecolor{Brick}{RGB}{183,48,46}
\definecolor{Amber}{RGB}{180,112,0}
\definecolor{Slate}{RGB}{54,65,82}
\definecolor{Steel}{RGB}{74,89,116}
\definecolor{Forest}{RGB}{32,120,80}
\definecolor{Ink}{RGB}{22,32,46}
\definecolor{PaperBlue}{RGB}{235,244,255}
\definecolor{PaperTeal}{RGB}{232,249,242}
\definecolor{PaperRed}{RGB}{255,239,239}
\definecolor{PaperAmber}{RGB}{255,247,232}
\definecolor{PaperGray}{RGB}{245,247,250}

\newcommand{\blue}[1]{\textcolor{Navy}{#1}}
\newcommand{\red}[1]{\textcolor{Brick}{#1}}
\newcommand{\green}[1]{\textcolor{Teal}{#1}}
\newcommand{\term}[1]{\textbf{\textcolor{Navy}{#1}}}
\newcommand{\warn}[1]{\textbf{\textcolor{Brick}{#1}}}
\newcommand{\focus}[1]{\textbf{\textit{\textcolor{Teal}{#1}}}}
\newcommand{\keysent}[1]{\textbf{\textit{\uline{#1}}}}
\let\origtextbf\textbf
\renewcommand{\textbf}[1]{\origtextbf{\textit{#1}}}

\hypersetup{
  colorlinks=true,
  linkcolor=Steel,
  urlcolor=Teal
}

%%%%%%%%%%%%%%%%%%%%%%%%%%%%%%%%%%%%%%%%
% SECTION STYLING
%%%%%%%%%%%%%%%%%%%%%%%%%%%%%%%%%%%%%%%%
\titleformat{\section}
{\LARGE\bfseries\color{Navy}}
{\thesection}{1em}{}

\titleformat{\subsection}
{\Large\bfseries\color{Teal}}
{\thesubsection}{1em}{}

\titleformat{\subsubsection}
{\large\bfseries\itshape\color{Steel}}
{\thesubsubsection}{1em}{}

\titlespacing*{\section}{0pt}{1.2em}{0.6em}
\titlespacing*{\subsection}{0pt}{0.9em}{0.4em}
\titlespacing*{\subsubsection}{0pt}{0.7em}{0.25em}

%%%%%%%%%%%%%%%%%%%%%%%%%%%%%%%%%%%%%%%%
% BOXES
%%%%%%%%%%%%%%%%%%%%%%%%%%%%%%%%%%%%%%%%
\newtcolorbox{bluebox}[1]{
breakable,
colback=PaperBlue,
colframe=Navy,
coltitle=white,
colbacktitle=Navy,
arc=2mm,
boxrule=1.0pt,
left=1.4mm,right=1.4mm,top=1.2mm,bottom=1.2mm,
title=#1,
fonttitle=\bfseries\large
}

\newtcolorbox{redbox}[1]{
breakable,
colback=PaperRed,
colframe=Brick,
coltitle=white,
colbacktitle=Brick,
arc=2mm,
boxrule=1.0pt,
left=1.4mm,right=1.4mm,top=1.2mm,bottom=1.2mm,
title=#1,
fonttitle=\bfseries\large
}

\newtcolorbox{greenbox}[1]{
breakable,
colback=PaperTeal,
colframe=Teal,
coltitle=white,
colbacktitle=Teal,
arc=2mm,
boxrule=1.0pt,
left=1.4mm,right=1.4mm,top=1.2mm,bottom=1.2mm,
title=#1,
fonttitle=\bfseries\large
}

\renewcommand{\arraystretch}{1.28}
\setlength{\tabcolsep}{6pt}
\arrayrulecolor{Slate}

%%%%%%%%%%%%%%%%%%%%%%%%%%%%%%%%%%%%%%%%
% GLOBAL READABILITY SETTINGS
%%%%%%%%%%%%%%%%%%%%%%%%%%%%%%%%%%%%%%%%
\setlist[itemize]{leftmargin=1.5em,itemsep=0.22em,topsep=0.35em,label=\textcolor{Teal}{\textbullet}}
\setlist[enumerate]{leftmargin=1.6em,itemsep=0.26em,topsep=0.35em}
\setlength{\parskip}{0.18em}
\setlength{\parindent}{0pt}
\linespread{1.08}

\lstset{
  basicstyle=\ttfamily\small\color{Ink},
  keywordstyle=\bfseries\color{Navy},
  commentstyle=\itshape\color{Forest},
  stringstyle=\color{Brick},
  numberstyle=\scriptsize\color{Steel},
  numbers=left,
  numbersep=7pt,
  frame=single,
  rulecolor=\color{Steel},
  framerule=0.5pt,
  backgroundcolor=\color{PaperGray},
  breaklines=true,
  showstringspaces=false,
  tabsize=2,
  xleftmargin=0.4em,
  xrightmargin=0.4em
}

%%%%%%%%%%%%%%%%%%%%%%%%%%%%%%%%%%%%%%%%
% DOC HEADER
%%%%%%%%%%%%%%%%%%%%%%%%%%%%%%%%%%%%%%%%
\title{\blue{NVIDIA AI/ML Infrastructure -- Study Review}}
\author{}
\date{}

\begin{document}
\color{Ink}
\maketitle

\tableofcontents
\newpage

% ======================================================
\section{Mental Model of AI/ML Infrastructure}
% ======================================================

\begin{bluebox}{Layered Stack}
Hardware $\rightarrow$ GPU Architecture $\rightarrow$ Network $\rightarrow$ Storage $\rightarrow$ Scheduler $\rightarrow$ Runtime $\rightarrow$ Framework $\rightarrow$ Model $\rightarrow$ Experiment Tracking
\end{bluebox}

\begin{greenbox}{Section 1 at a Glance (Recall Outline)}
\begin{enumerate}[leftmargin=1.3em]
\item \textbf{Hardware:} How cluster and node topology is wired
\item \textbf{GPU Architecture:} SM, warp, memory hierarchy, precision units
\item \textbf{Network:} Ethernet vs InfiniBand, RDMA/GPUDirect, management planes
\item \textbf{Storage:} Fast-to-slow hierarchy and I/O bottleneck signs
\item \textbf{Scheduler:} Slurm/Kubernetes resource orchestration
\item \textbf{Runtime/Framework:} CUDA + NCCL + PyTorch distributed stack
\item \textbf{Model/Tracking:} Model shape drives infra; tracking validates impact
\end{enumerate}
\end{greenbox}

\begin{bluebox}{One-Line Exam Mental Model}
\textbf{If GPUs are idle, find the starvation layer:} input (\textbf{storage/data}), communication (\textbf{network/NCCL}), or runtime/configuration (\textbf{scheduler/framework}).
\end{bluebox}

\subsection{Hardware}

\begin{bluebox}{Hardware Recall}
\textbf{Path:} User $\rightarrow$ Login Node $\rightarrow$ Compute Node(s) $\leftrightarrow$ Shared Storage
\end{bluebox}

\textbf{HPC stack view (cluster-level):}
\begin{itemize}
\item User connects to a login/head node.
\item Login/head nodes connect to compute nodes.
\item Compute nodes share a parallel file system and are interconnected.
\item Storage systems are typically separate from compute clusters.
\end{itemize}

\textbf{Inside a compute node:}
\begin{itemize}
\item Main components: CPU, GPU, and (in many modern systems) DPU/SmartNIC.
\item CPU and GPU are connected through PCIe.
\item GPUs within a node are connected through NVLink or NVSwitch.
\end{itemize}

\subsection{GPU Architecture}

\begin{bluebox}{GPU Recall}
\textbf{Think in levels:} GPU $\rightarrow$ SM $\rightarrow$ Warp (32 threads) $\rightarrow$ Execution units + memory hierarchy.
\end{bluebox}

\begin{itemize}
\item Each GPU has multiple \textbf{SMs} (Streaming Multiprocessors).
\item SMs share global memory (HBM/DRAM) through shared L2.
\item Inside each SM: register files, warp scheduler, dispatch logic, L1/shared memory, and execution cores.
\item A warp is typically 32 threads, and schedulers pick ready warps to improve occupancy and hide latency.
\item Cores may target different precisions: Tensor Core ops, FP32, INT, etc.
\end{itemize}

\subsection{Network}

\begin{redbox}{Bottleneck Clue}
If network or storage is busy while GPUs are idle, training is often stalled by gradient synchronization or data movement.
\end{redbox}

\begin{bluebox}{Network Recall}
\textbf{Data-plane speed tools:} InfiniBand + RDMA (+ GPUDirect where supported).  
\textbf{Control/ops planes:} in-band for normal ops, out-of-band for hardware recovery/management.
\end{bluebox}

\begin{itemize}
\item Ethernet is commonly used for general cluster traffic.
\item InfiniBand is commonly used as the high-speed training interconnect.
\item RDMA enables low-overhead data movement with minimal CPU involvement.
\item GPUDirect RDMA/GPUDirect Storage can reduce unnecessary host copies in supported setups.
\item DMA: direct memory transfer between device and host memory without active CPU copying loops.
\item In-band network: regular OS/service traffic (SSH, scheduler, etc.).
\item Out-of-band management network: hardware management path even if host OS is down.
\end{itemize}

\subsection{Storage}

\begin{redbox}{Bottleneck Clue}
If DataLoader is maxed out and GPUs are idle, storage/input pipeline is likely the bottleneck.
\end{redbox}

\begin{bluebox}{Storage Recall}
\textbf{Order (fast to slow):} NVMe $\rightarrow$ Parallel FS $\rightarrow$ NFS $\rightarrow$ Object Store.
\end{bluebox}

\textbf{Fastest to slowest (practical training view):}
\begin{enumerate}[leftmargin=1.4em]
\item \textbf{Local NVMe}: fastest scratch storage for temporary datasets/intermediate files.
\item \textbf{Parallel file systems} (Lustre, GPFS, BeeGFS): shared high-throughput storage across nodes.
\item \textbf{NFS}: good for configs/scripts/small files, not ideal for heavy training I/O.
\item \textbf{Object storage} (S3-style): good for logs/checkpoints/archive, typically higher latency for hot training reads.
\end{enumerate}

\subsection{Scheduler}

\textbf{Common options:} Slurm or Kubernetes.

\begin{bluebox}{Scheduler Recall}
\textbf{Slurm:} batch/HPC-first job scheduling.  
\textbf{Kubernetes:} container-first orchestration with GPU plugins.
\end{bluebox}

{\scriptsize
\setlength{\tabcolsep}{2.6pt}
\begin{tabularx}{0.98\textwidth}{|>{\raggedright\arraybackslash}p{1.5cm}|>{\raggedright\arraybackslash}X|>{\raggedright\arraybackslash}X|}
\hline
\textbf{Area} & \textbf{Slurm (high-level)} & \textbf{Kubernetes (high-level)} \\
\hline
Workload model & Batch/HPC scheduler for queued jobs and resource allocation. & Container orchestrator for pod-based workloads. \\
\hline
GPU resource model & GPU resources configured via \texttt{slurm.conf} and \texttt{gres.conf}. & GPUs requested through Kubernetes resource specs and exposed via device plugins (e.g., NVIDIA plugin). \\
\hline
\shortstack[l]{Scheduling\\fairness} & Priorities, QoS, fairshare, and preemption policies. & Scheduler places pods on nodes with available GPUs; policies are cluster-config driven. \\
\hline
\shortstack[l]{Monitoring\\ops} & \texttt{sacct}, \texttt{sreport}, dashboards, and node/GPU health logs. & Strong observability ecosystem (e.g., Prometheus/Grafana), plus autoscaling support. \\
\hline
\shortstack[l]{Performance\\capacity} & Tune GPU placement, NCCL config, CPU/GPU affinity, and network settings; plan capacity from workload trends. & Optimize placement and scaling behavior through scheduling constraints and autoscaling signals. \\
\hline
\end{tabularx}
}

\vspace{0.4cm}
\textbf{Slurm example command:}
\begin{center}
\small
\texttt{srun --partition=gpu --gres=gpu:1 --mem=16G}\\
\texttt{--time=01:00:00 --pty bash}
\normalsize
\end{center}

\begin{itemize}
\item \texttt{srun}: starts an interactive run.
\item \texttt{--gres=gpu:1}: requests one GPU.
\item \texttt{--mem=16G --time=01:00:00}: sets memory and time limits.
\item \texttt{--pty bash}: opens an interactive shell on the allocated node.
\end{itemize}

\subsection{Runtime}
\begin{bluebox}{Runtime Recall}
\textbf{Runtime stack:} CUDA kernels + NCCL collectives + drivers + container environment.
\end{bluebox}

\begin{itemize}
\item CUDA runtime and kernels (compiled with \texttt{nvcc}).
\item NCCL for distributed GPU communication.
\item Containers for reproducibility.
\item Drivers as the host-to-GPU interface.
\end{itemize}

Common issues: driver/runtime mismatch, inefficient kernel launches, communication backend bottlenecks, or container misconfiguration.

\subsection{Framework}
\textbf{PyTorch}

\begin{bluebox}{Framework Recall}
\textbf{DDP:} replicate model, sync gradients.  
\textbf{FSDP:} shard model states to fit larger models.  
\textbf{TP/PP:} split model compute across dimensions/stages.
\end{bluebox}

\textbf{DDP (DistributedDataParallel):}
\begin{itemize}
\item One model replica per GPU; gradients synchronized via AllReduce.
\item Typical setup: wrap model with DDP, use \texttt{DistributedSampler}, call \texttt{loss.backward()} normally.
\item Scaling can suffer with very small per-GPU batches because communication dominates.
\end{itemize}

\textbf{FSDP (Fully Sharded Data Parallel):}
\begin{itemize}
\item Shards parameters, gradients, and optimizer states across ranks.
\item Useful for fitting larger models.
\item Combine with mixed precision/checkpointing for memory efficiency.
\end{itemize}

\textbf{Tensor and pipeline parallelism:}
\begin{itemize}
\item Key when model size exceeds what plain data parallelism can handle.
\item Directly affects communication pattern, memory pressure, and scaling behavior.
\end{itemize}

\subsection{Model}
\begin{bluebox}{Model Recall}
Model architecture determines memory pressure, communication pattern, and best parallelism strategy.
\end{bluebox}

Model size, structure, and compute pattern heavily influence infrastructure decisions: parallelism strategy, memory requirements, communication overhead, and inference latency.

\subsection{Experiment Tracking}
\begin{bluebox}{Tracking Recall}
No tracking means no reliable diagnosis, no reproducible comparisons, and weak optimization decisions.
\end{bluebox}

Tools like Weights \& Biases and monitoring dashboards provide visibility into training performance and regressions. Without tracking, it is difficult to diagnose failures, evaluate optimizations, or scale workflows reliably.

\section{PyTorch Distributed Components by Responsibility}

{\scriptsize
\setlength{\tabcolsep}{3pt}
\begin{tabularx}{0.98\textwidth}{|>{\raggedright\arraybackslash}p{2.2cm}|>{\raggedright\arraybackslash}p{2.3cm}|>{\raggedright\arraybackslash}p{2.5cm}|>{\raggedright\arraybackslash}X|}
\hline
\textbf{Responsibility} & \textbf{Component} & \textbf{What It Does} & \textbf{Implementation Steps} \\
\hline
Multi-process launching & Torchrun & Spawns one process per GPU and sets rank/world size & Launch with \texttt{torchrun --nproc\_per\_node=N}; read \texttt{LOCAL\_RANK}; set device; init process group. \\
\hline
Process group init & init\_process\_group & Creates communication group (NCCL/Gloo) & Call \texttt{torch.distributed.init\_process\_group(...)} early with \texttt{backend="nccl"} and set local CUDA device. \\
\hline
Distributed data parallel & DDP & Replicates model on each GPU and syncs gradients & Wrap model in DDP, use \texttt{DistributedSampler}, call \texttt{loss.backward()}. \\
\hline
Distributed sampling & DistributedSampler & Shards dataset across ranks & Pass sampler into DataLoader and call \texttt{sampler.set\_epoch(epoch)}. \\
\hline
Fully sharded data parallel & FSDP & Shards parameters, gradients, and optimizer state & Use torch.distributed.fsdp.FullyShardedDataParallel and configure sharding strategy, mixed precision, and wrapping policy. \\
\hline
Sharded checkpointing & dist.checkpoint & Saves/restores large distributed states efficiently & Use \texttt{torch.distributed.checkpoint} with FSDP-compatible state-dict APIs for save/load. \\
\hline
Gradient accumulation & no\_sync() (DDP) & Avoids sync on intermediate microsteps & Use \texttt{with model.no\_sync():} and sync on final accumulation step. \\
\hline
Mixed precision & autocast + GradScaler & Speed + memory improvements with lower precision & Wrap forward pass in \texttt{autocast()}, scale gradients with \texttt{GradScaler}. \\
\hline
Activation checkpointing & checkpoint API & Trades compute for lower memory usage & Use torch.utils.checkpoint on large blocks/layers and recompute on backward. \\
\hline
Tensor parallelism & device\_mesh + DTensor & Splits tensors/layers across GPUs & Define mesh, shard tensors, coordinate collectives. \\
\hline
Pipeline parallelism & Pipeline APIs/ecosystem tools & Splits model into stages via microbatches & Partition model stages and tune schedule to reduce bubbles. \\
\hline
Communication backend & NCCL & Handles GPU collectives & Use \texttt{backend="nccl"} and tune/debug via environment settings. \\
\hline
Compilation / optimization & torch.compile & Graph capture and kernel fusion & Compile model and validate with profiler-based speedups. \\
\hline
Data loading optimization & DataLoader tuning & Keeps GPUs fed & Tune \texttt{num\_workers}, \texttt{pin\_memory}, prefetch settings. \\
\hline
Profiling & torch.profiler & Captures CPU/GPU/comm timelines & Run scheduled traces and inspect bottlenecks. \\
\hline
Fault tolerance / elastic & TorchElastic (\texttt{torchrun}) & Supports restarts and membership changes & Use rendezvous backend + robust checkpoint resume flow. \\
\hline
\end{tabularx}
}

\subsection{Explain End-to-End Training Flow}
Whiteboard explanation from data loading to optimizer step across 32 GPUs.

\vspace{0.5cm}

\begin{greenbox}{Likely Question}
Walk me through how a distributed PyTorch job runs across multiple nodes.
\end{greenbox}

Typical flow:
\begin{itemize}
\item Start with a launcher (\texttt{torchrun}, often orchestrated by Slurm/Kubernetes).
\item Spawn one process per GPU and assign each process a global rank, local rank, and world size.
\item Each process sets its CUDA device and initializes the distributed process group (usually NCCL).
\item Model is wrapped in DDP (or FSDP for sharded training).
\item Dataset is partitioned with \texttt{DistributedSampler} so each rank gets a different data shard.
\item During training, each rank runs forward/backward locally.
\item Gradient synchronization (AllReduce or sharded collectives) keeps model updates consistent.
\item Optimizer step proceeds in lockstep across ranks.
\end{itemize}

Overall performance depends on compute, network bandwidth, storage throughput, data loader efficiency, and scheduler configuration. A bottleneck in any layer can reduce scaling efficiency.

% ======================================================
\section{Distributed Training -- FULL DEPTH}
% ======================================================

\subsection{DDP Internals}

A mini-batch is a small subset of training examples processed in one forward + backward pass.

\begin{bluebox}{DDP Mental Model (Forward to Optimizer Step)}
\begin{enumerate}[leftmargin=1.3em]
\item Forward pass runs independently on each GPU replica.
\item Backward pass computes parameter gradients.
\item Autograd hooks trigger communication when gradient buckets are ready.
\item NCCL AllReduce synchronizes gradients across ranks.
\item Optimizer step updates weights in lockstep across GPUs.
\end{enumerate}
\end{bluebox}

\begin{bluebox}{What Triggers Communication in DDP}
\begin{itemize}
\item DDP registers autograd hooks on parameters.
\item Hooks are attached to gradient accumulation logic.
\item As gradients are produced in backward, they are added to buckets.
\item When a bucket is ready, DDP schedules NCCL collective communication.
\end{itemize}
\end{bluebox}

\begin{redbox}{Why Overlap Matters}
Communication starts during backward (not only after backward ends).  
This overlap of backward compute + NCCL communication reduces GPU idle time and hides communication latency.
\end{redbox}

\textbf{Straggler effect (preserved + organized):}
\begin{itemize}
\item Faster GPUs wait for slower ones at synchronization points.
\item Typical causes: data imbalance, CPU data-loading bottlenecks, storage I/O stalls, or network congestion.
\item Quick data check: compare local batch sizes/runtimes across ranks (for example, print/trace per-rank batch stats).
\end{itemize}

\textbf{Debug heuristics (preserved + corrected):}
\begin{itemize}
\item \textbf{Data imbalance:} one rank receives more work; that GPU is idle \textbf{before} forward starts on next step.
\item \textbf{CPU bottleneck:} DataLoader workers too few; CPU saturated (check with \texttt{htop}).
\item \textbf{Disk/storage bottleneck:} slow reads from storage or network filesystem (check with \texttt{iostat} and storage metrics).
\item \textbf{Network bottleneck:} GPUs wait during gradient sync, usually \textbf{after} backward kernels (check NCCL traces/\texttt{nsys}; use network tools such as \texttt{ifstat}/fabric metrics).
\item \textbf{GPU utilization clues:} \texttt{nvidia-smi} shows low utilization or long idle gaps on specific ranks.
\end{itemize}


Topics to master:
\begin{itemize}
\item Autograd hooks
\item Gradient bucketing: AllReduce starts when a gradient bucket is full, which enables overlap with backward compute.
\item Bucket-size tradeoff: smaller buckets improve overlap but increase communication-call overhead; larger buckets reduce call count but can delay communication start.
\item Tuning bucket size can improve scaling by balancing overlap vs communication overhead.
\item Overlapping compute + communication
\item NCCL backend: optimized for GPU collectives, uses NVLink/PCIe topology awareness, and builds efficient communication rings/trees.
\item Rank initialization
\end{itemize}

\begin{redbox}{Follow-Up They Might Ask}
Why does small batch size hurt DDP scaling?
\begin{itemize}
\item With very small per-GPU batch size, compute per step is short while communication cost stays similar, so communication dominates.
\item Increasing global batch size often requires learning-rate retuning (and sometimes warmup adjustments) for stable convergence.
\end{itemize}
\end{redbox}


At the end, overlap means backward compute continues while previous gradients are being reduced.

\textbf{How to detect poor overlap:}
\begin{itemize}
\item Large communication-only gaps after backward.
\item High GPU idle time near the step tail.
\item In profiler timelines (e.g., Nsight Systems), good runs show interleaving of backward kernels and NCCL ops; poor runs show big NCCL blocks after backward completion.
\end{itemize}


\begin{bluebox}{DDP Init and Failure Checklist}
Each DDP process has a unique \texttt{RANK}; \texttt{WORLD\_SIZE} is total processes.

\textbf{Common env vars:}
\begin{itemize}
\item \texttt{MASTER\_ADDR}
\item \texttt{MASTER\_PORT}
\item \texttt{RANK}
\item \texttt{WORLD\_SIZE}
\end{itemize}

\textbf{If setup is wrong, typical symptoms:}
\begin{itemize}
\item rank mismatch or world-size mismatch \(\rightarrow\) hang/deadlock at init/collectives
\item wrong master address/port conflict/firewall block \(\rightarrow\) rendezvous failure
\item Slurm misconfiguration \(\rightarrow\) inconsistent rank mapping
\end{itemize}
\end{bluebox}

\begin{bluebox}{DDP Step-by-Step (Interview Version)}
Each GPU runs a full model replica on a different data shard. During backward, autograd hooks bucket gradients and trigger NCCL AllReduce. After synchronization, each rank applies the same optimizer step.  
\textbf{Key point:} gradient sync is overlapped during backward, not only after backward ends.
\end{bluebox}

\begin{bluebox}{PyTorch FSDP: What Changes vs Vanilla Training}
\begin{itemize}
\item launch one process per GPU (\texttt{torchrun} or Slurm launcher)
\item initialize distributed process group (\texttt{init\_process\_group})
\item set local CUDA device per rank
\item wrap model/submodules with FSDP
\item train with normal forward/backward/step loop
\item use FSDP-aware sharded checkpointing
\end{itemize}

\textbf{Key APIs to know:}
\begin{itemize}
\item \texttt{torchrun --nproc\_per\_node=...}
\item \texttt{torch.distributed.init\_process\_group}
\item \texttt{torch.distributed.fsdp.FullyShardedDataParallel}
\item \texttt{torch.distributed.fsdp.wrap} (auto-wrap policy)
\item \texttt{ShardingStrategy.FULL\_SHARD} (ZeRO-3-like behavior)
\item \texttt{FSDP.state\_dict\_type(..., SHARDED\_STATE\_DICT)}
\end{itemize}
\end{bluebox}

\begin{lstlisting}[language=Python, caption={FSDP pseudocode (PyTorch)}]
import os
import torch
import torch.distributed as dist
from torch.optim import AdamW
from torch.utils.data import DataLoader, DistributedSampler
from torch.distributed.fsdp import FullyShardedDataParallel as FSDP
from torch.distributed.fsdp import ShardingStrategy, StateDictType
from torch.distributed.fsdp.api import MixedPrecision
from torch.distributed.fsdp.wrap import size_based_auto_wrap_policy

def setup_dist():
    dist.init_process_group(backend="nccl")
    local_rank = int(os.environ["LOCAL_RANK"])
    torch.cuda.set_device(local_rank)
    return local_rank, dist.get_rank(), dist.get_world_size()

def wrap_with_fsdp(model):
    auto_wrap = lambda m, recurse, nonwrapped_numel: size_based_auto_wrap_policy(
        m, recurse=recurse, nonwrapped_numel=nonwrapped_numel, min_num_params=1e6
    )
    mp = MixedPrecision(
        param_dtype=torch.bfloat16,
        reduce_dtype=torch.bfloat16,
        buffer_dtype=torch.bfloat16,
    )
    return FSDP(
        model,
        auto_wrap_policy=auto_wrap,
        sharding_strategy=ShardingStrategy.FULL_SHARD,
        mixed_precision=mp,
        device_id=torch.cuda.current_device(),
    )
\end{lstlisting}

\begin{redbox}{FSDP FULL\_SHARD Recall}
FULL\_SHARD is ZeRO-3-like: parameters, gradients, and optimizer states are sharded.  
Parameters are gathered just-in-time per wrapped unit; gradients are reduce-scattered in backward.
\end{redbox}

\begin{bluebox}{DeepSpeed ZeRO: What Changes vs Vanilla Training}
\begin{itemize}
\item define DeepSpeed config with \texttt{zero\_optimization.stage}
\item call \texttt{deepspeed.initialize(...)}
\item train with \texttt{engine.backward(loss)} and \texttt{engine.step()}
\item save ZeRO-aware checkpoints via \texttt{engine.save\_checkpoint(...)}
\end{itemize}
\end{bluebox}

\begin{lstlisting}[language=Python, caption={DeepSpeed ZeRO pseudocode}]
import deepspeed
import torch
from torch.utils.data import DataLoader, DistributedSampler

def get_ds_config(stage: int):
    return {
        "train_batch_size": GLOBAL_BS,
        "train_micro_batch_size_per_gpu": PER_GPU_BS,
        "gradient_accumulation_steps": GRAD_ACCUM,
        "bf16": {"enabled": True},
        "zero_optimization": {"stage": stage},
        "optimizer": {
            "type": "AdamW",
            "params": {"lr": LR, "betas": [0.9, 0.999], "eps": 1e-8, "weight_decay": WD}
        },
    }
\end{lstlisting}

\begin{bluebox}{Quick Comms Notes}
\textbf{Ring AllReduce:} gradients are split and passed around a ring; each rank talks to neighbors only, which is bandwidth-efficient.  
\textbf{Gradient accumulation:} accumulate multiple micro-batches before optimizer step to reduce sync frequency.
\end{bluebox}

\begin{redbox}{Scaling/Debug Interview Notes}
\textbf{Why efficiency drops at high GPU count:} communication/synchronization dominates compute, latency and stragglers hurt more.  
\textbf{Common causes:} data-loader bottlenecks, too-small per-GPU batch, communication blocking, I/O saturation, CPU under-provisioning.

\textbf{Job hang checklist:}
\begin{itemize}
\item rank/world-size mismatch
\item one crashed rank
\item NCCL deadlock or collective-order mismatch
\item firewall/port/master-address issues
\item Slurm rank mapping errors
\end{itemize}

\textbf{Single-node vs multi-node:}
\begin{itemize}
\item single node (NVLink): higher bandwidth, lower latency, fewer failure points
\item multi-node (InfiniBand/Ethernet): higher latency, more sync overhead, topology matters
\end{itemize}
\end{redbox}

\vspace{1cm}

\subsection{FSDP / ZeRO -- Memory Deep Dive}

\begin{bluebox}{Baseline Memory in Vanilla DDP}
Each GPU stores full parameters + full gradients + full optimizer states.  
Per-GPU memory scales poorly for very large models because model states are fully replicated.
\end{bluebox}

ZeRO stands for \textbf{Zero Redundancy Optimizer}. It removes redundant model-state copies by sharding them across GPUs. FSDP is PyTorch's ZeRO-style implementation.

Understand:
\begin{itemize}
\item \textbf{Parameter sharding:} each GPU stores only \(1/N\) parameters at rest; parameters are gathered just in time per layer.
\item \textbf{Gradient partitioning:} gradients are partitioned (typically via ReduceScatter), so each rank keeps only its shard.
\item \textbf{Optimizer state partitioning:} major memory win; for Adam, optimizer state is often the largest memory component.
\item \textbf{Activation checkpointing:} save fewer activations and recompute in backward (less memory, more compute).
\item \textbf{Communication amplification:} ZeRO-3/FSDP reduces memory but increases communication events (gather/scatter).
\end{itemize}

\begin{greenbox}{Prompt}
Explain tradeoffs between ZeRO-2 and ZeRO-3.

\textbf{ZeRO-2}
\begin{itemize}
\item shards optimizer states and gradients
\item parameters remain replicated
\item lower communication overhead than ZeRO-3
\item easier to scale when network bandwidth is limited
\end{itemize}

\textbf{ZeRO-3}
\begin{itemize}
\item shards parameters + gradients + optimizer states
\item maximum memory savings (roughly model-state memory \(\sim 1/N\) per GPU)
\item higher communication and more synchronization points
\item enables larger model sizes at the cost of communication complexity
\end{itemize}

\textbf{Memory sketch:}
\begin{itemize}
\item DDP: \(\text{params} + \text{grads} + \text{optimizer}\)
\item ZeRO-2: \(\text{params} + \frac{\text{grads}}{N} + \frac{\text{optimizer}}{N}\)
\item ZeRO-3: \(\frac{\text{params}}{N} + \frac{\text{grads}}{N} + \frac{\text{optimizer}}{N}\)
\end{itemize}

\textbf{When are parameters gathered in FSDP?}
\begin{itemize}
\item just before a wrapped layer runs forward
\item can be freed after layer compute
\item gathered layer-by-layer (not entire model at once)
\end{itemize}

\textbf{ReduceScatter vs AllReduce}
\begin{itemize}
\item AllReduce: all ranks end with full reduced gradients
\item ReduceScatter: reduce + shard in one step; each rank keeps only its shard
\end{itemize}
\end{greenbox}

\begin{bluebox}{Memory-Spike Debug Guide (FSDP/ZeRO)}
\textbf{Where to inspect:} forward gather, backward reduce-scatter, checkpointing, accumulation windows.  
Use \texttt{torch.cuda.memory\_summary()}, Nsight Systems, and PyTorch profiler.

\textbf{Common spike causes:}
\begin{itemize}
\item coarse wrapping granularity (gathers too much at once)
\item forward prefetch overlaps too many large gathers
\item consolidated checkpoints gathering full model on rank 0
\item mixed-precision misconfiguration creating unexpected FP32 copies
\item large gradient-accumulation factor increasing peak live tensors
\end{itemize}

\textbf{Checkpointing note:}
\begin{itemize}
\item sharded checkpoints are most scalable for ZeRO/FSDP
\item full consolidated checkpoints are easier to consume but can cause rank-0 memory spikes
\end{itemize}
\end{bluebox}

\vspace{2cm}

% ======================================================
\section{Parallelism Strategies for LLMs}
% ======================================================
\begin{bluebox}{Memory Reminder}
What usually blows up memory:
\begin{itemize}
\item steady-state memory: parameters, gradients, optimizer states
\item peak memory: activations saved for backward (often grows with batch size and sequence length)
\end{itemize}
\end{bluebox}

\subsection{Data Parallelism}
\begin{itemize}
\item Input batch is split across devices/processes.
\item Each device runs forward + backward on its mini-batch independently.
\item Pro: simple scaling, no pipeline bubbles, independent execution.
\item Con: each device stores a full copy of model states (parameters + gradients + optimizer), so memory footprint is not reduced.
\item Gradient synchronization happens through collectives (commonly AllReduce or ReduceScatter + AllGather, depending on implementation).
\end{itemize}

\subsection{Tensor Parallelism}
\begin{bluebox}{Tensor Parallelism (Intra-Layer)}
Split tensors within a layer (for example, sharded matmul weights). This reduces per-device compute and per-device parameter footprint for that layer, but increases inter-device communication.
\end{bluebox}

\textbf{Row-wise vs column-wise split:}
\begin{itemize}
\item \textbf{Column-parallel linear (split output features):} each GPU computes a slice of output; next layers may need AllGather (or compatible sharding). Backward usually includes gradient communication.
\item \textbf{Row-parallel linear (split input features):} each GPU consumes input slices and partial outputs are summed (commonly via AllReduce).
\item Frameworks place AllGather/AllReduce where cost is lowest.
\end{itemize}

\begin{itemize}
\item \textbf{Splitting weight matrices}
\item Heavy Transformer ops affected: attention projections (Q/K/V/output) and FFN linear layers.
\item Example (column split): for \(W \in \mathbb{R}^{d \times 4d}\), shard columns across GPUs (e.g., GPU0: \(W[:,0:2d]\), GPU1: \(W[:,2d:4d]\)); each GPU computes local \(XW_i\), then combine as needed.
\item \textbf{AllGather requirements}
\item Partial outputs often must be gathered before the next layer.
\item Gradients are reduced in backward (common patterns: AllGather, ReduceScatter, AllReduce).
\item This communication occurs repeatedly across layers.
\item \textbf{Communication cost}
\item Tensor parallelism introduces per-layer communication across TP groups.
\item Communication cost scales with both layer count and TP size.
\end{itemize}

\subsection{Pipeline Parallelism}
\begin{bluebox}{Pipeline Parallelism}
Layer-wise model split across devices. Memory is reduced per device because each stage stores only its own layers, but scheduling is required to minimize idle time (pipeline bubbles).
\end{bluebox}

\begin{itemize}
\item Microbatching
\item Without microbatching, devices idle heavily across stages.
\item With microbatching, the pipeline fills:
\begin{itemize}
\item Time 1: GPU0 \(\rightarrow\) MB1
\item Time 2: GPU0 \(\rightarrow\) MB2, GPU1 \(\rightarrow\) MB1
\item Time 3: GPU0 \(\rightarrow\) MB3, GPU1 \(\rightarrow\) MB2, GPU2 \(\rightarrow\) MB1
\end{itemize}
\item Bubble overhead
\item Idle time at start/end is the pipeline bubble.
\item Bubble cost reminder: \((\#\text{stages} - 1)/\#\text{microbatches}\).
\item More microbatches generally means a smaller bubble.
\item Stage balancing
\item If one stage is heavier, the whole pipeline stalls.
\item Stages should be balanced for both compute and memory.
\end{itemize}

\subsection{Hybrid Parallelism}
Explain how large models combine DP (across replicas) + TP (within nodes) + PP (across layers).

\vspace{2cm}

\begin{greenbox}{Design Prompt}
Design distributed training for a 100B parameter model on 64 GPUs.
\textbf{Placeholder retained:} significance of this exact number to be filled later.

Why combine DP + TP + PP:
\begin{itemize}
\item TP helps when a single layer is too large for one GPU.
\item PP helps when the full model is too large for one GPU.
\item DP scales throughput across replica groups.
\end{itemize}

Communication differences:
\begin{itemize}
\item TP: heavy per-layer communication.
\item PP: activation passing between stages.
\item DP: gradient synchronization.
\item Each type introduces a different bottleneck.
\end{itemize}
\end{greenbox}

\textbf{Questions to expect:}
\begin{itemize}
\item How would you train a 100B model across 64 GPUs?
\item What limits tensor-parallel scaling?
\item What limits pipeline scaling?
\item When do you prefer TP over PP?
\item When do you prefer PP over TP?
\end{itemize}

\textbf{Answer cues from your notes:}
\begin{itemize}
\item combine TP (split large layers), PP (partition layers across stages), and DP (scale throughput across groups)
\item TP limits: per-layer communication, interconnect bandwidth, latency
\item PP limits: bubble overhead, stage imbalance, activation memory
\item prefer TP when layers are extremely large and fast interconnect (e.g., NVLink) is available
\item prefer PP when the full model does not fit a single GPU and cross-node bandwidth constraints dominate
\end{itemize}

% ======================================================
\section{NCCL + Networking}
% ======================================================

\begin{bluebox}{NCCL Recall}
NCCL provides distributed collective communication primitives for multi-GPU training.
\end{bluebox}

\textbf{Core collectives from your notes:}
\begin{itemize}
\item \textbf{Broadcast:} one rank sends data to all ranks (often weights/config).
\item \textbf{AllReduce:} reduce across all ranks and store result on every rank (gradient sync).
\item \textbf{Reduce:} reduce across ranks and store result on one rank (aggregation).
\item \textbf{AllGather:} gather shards so each rank has full data (common in unshard steps, e.g., FSDP).
\item \textbf{ReduceScatter:} reduce + scatter so each rank keeps only a shard (common for sharded gradients).
\end{itemize}

\textbf{Gradient-sync reminder:}
In data parallel training, each GPU computes partial gradients on its mini-batch, then synchronization collectives are used before weight update.

\subsection{AllReduce Algorithms}
AllReduce combines tensors across GPUs and distributes reduced results to all ranks. It is used to synchronize gradients and keep model state consistent.

\begin{itemize}
\item Ring
\item GPUs form a logical ring and exchange chunks with neighbors.
\item Communication pattern is \(2 \times (N-1)\) steps.
\item Tensor is split into chunks; reduce-scatter style pass plus all-gather pass.
\item Pro: bandwidth-efficient, scales well for large messages.
\item Con: latency accumulates with larger \(N\); weaker for small tensors.
\item Interview cue: if one GPU is slower, the ring stalls (straggler effect).
\item Tree
\item Pair GPUs in a reduction tree, combine upward, then broadcast down.
\item Complexity intuition: \(O(\log N)\).
\item Pro: lower latency, good for small messages.
\item Con: less bandwidth-efficient; root/intermediate nodes can become hot spots.
\item Prefer tree when latency dominates.
\item Hierarchical
\item Common in multi-node setups.
\item Two-stage reduction:
\begin{itemize}
\item intra-node reduction (e.g., NVLink domain)
\item inter-node reduction (e.g., InfiniBand/Ethernet fabric)
\end{itemize}
\item Reduces cross-node traffic by exploiting faster intra-node links first.
\end{itemize}

\subsection{InfiniBand + RDMA}
\begin{itemize}
\item What is RDMA?
\item Direct memory access across machines with minimal CPU/kernel involvement in the data path.
\item Path intuition from your notes:
\begin{itemize}
\item traditional: GPU \(\rightarrow\) CPU/kernel \(\rightarrow\) NIC \(\rightarrow\) network \(\rightarrow\) NIC/kernel \(\rightarrow\) CPU \(\rightarrow\) GPU
\item RDMA-style path reduces host-copy overhead
\end{itemize}
\item How does EFA work?
\item EFA = Elastic Fabric Adapter.
\item Enables low-latency networking in cloud HPC/distributed workloads.
\item Supports high-speed NCCL collective communication patterns in cloud settings.
\item Why bypass CPU?
\item CPU context switching, kernel interrupts, and extra copies increase latency and overhead.
\item Latency vs bandwidth
\item Large gradient tensors are often bandwidth-critical.
\item Small control messages are latency-critical.
\item Scaling behavior depends on both.
\end{itemize}

\begin{redbox}{Troubleshooting Scenario}
Scaling from 8 $\rightarrow$ 32 GPUs drops efficiency from 85\% to 60\%. Diagnose.

\textbf{Diagnosis flow (from your notes):}
\begin{enumerate}[leftmargin=1.3em]
\item Check GPU utilization first (\texttt{nvidia-smi}).
\begin{itemize}
\item 90--100\%: likely compute-bound
\item 40--60\%: often communication or I/O bound
\item very low utilization: likely stalls
\end{itemize}
\item Profile timelines (\texttt{nsys}) and verify if NCCL dominates.
\item Use PyTorch profiler to measure NCCL time vs compute-kernel time.
\item Compute scaling efficiency and flag low values (your note: \(<70\%\) suspicious).
\end{enumerate}

\textbf{Network and communication checks:}
\begin{itemize}
\item Enable NCCL debug logs:
\begin{itemize}
\item \texttt{export NCCL\_DEBUG=INFO}
\item \texttt{export NCCL\_DEBUG\_SUBSYS=ALL}
\end{itemize}
\item Inspect InfiniBand/network tools: \texttt{ibstat}, \texttt{ibv\_devinfo}, \texttt{sar -n DEV 1}, \texttt{iftop}, \texttt{nload}
\item Check topology and PCIe/NVLink paths: \texttt{nvidia-smi topo -m}
\item Watch for stragglers: one rank consistently slower/lower utilization
\item Check CPU loader pressure (\texttt{htop}), NUMA placement, and thermal throttling (\texttt{nvidia-smi -q -d TEMPERATURE})
\end{itemize}

\textbf{Batch-size effect from your notes:}
\[
\text{per-GPU batch} = \frac{\text{global batch}}{\text{\#GPUs}}
\]
If GPU count grows while global batch is fixed, per-GPU batch shrinks and communication can dominate.  
Typical fix in your notes: increase global batch and/or use gradient accumulation.

\textbf{Quick reference flow:}
\begin{itemize}
\item check GPU utilization
\item profile NCCL time
\item enable NCCL logs
\item check network saturation
\item check per-GPU batch size
\item look for stragglers
\item confirm hierarchical reduction behavior
\end{itemize}

\textbf{Strong answer cues for ``improve 32-GPU scaling'':}
\begin{itemize}
\item increase batch size
\item tune gradient bucket size
\item enable/tune hierarchical AllReduce path
\item reduce synchronization frequency
\item improve network
\end{itemize}
\end{redbox}

\vspace{2cm}

% ======================================================
% ======================================================
\section{GPU Architecture + Performance}
% ======================================================

\begin{bluebox}{Performance Reminder}
\[
\text{TFLOPs} \propto \text{cores} \times \text{clock} \times \text{FLOPs/cycle}
\]
Performance is strongly affected by memory bandwidth, available TFLOPs, and Tensor Core utilization.
\end{bluebox}

\subsection{SM Architecture}
\begin{greenbox}{SM Quick Recall}
\begin{itemize}
\item \textbf{Warp:} 32 threads per warp (SIMT style).
\item \textbf{Warp scheduler:} picks ready warps to hide latency.
\item \textbf{Dispatcher:} issues instructions to pipelines/cores.
\item If one warp stalls on memory, the scheduler switches to another warp.
\item If threads in one warp diverge on control flow, paths serialize and throughput drops.
\item \textbf{Occupancy:} active warps / max possible warps per SM.
\item Achieved occupancy can differ from theoretical occupancy due to scheduling/resource limits.
\item \textbf{Shared memory:} enables intra-block reuse and reduced global-memory traffic.
\item \textbf{Registers:} fastest storage tier for thread-local values.
\end{itemize}
\end{greenbox}

\keysent{Library/tool notes from your content:}\par
{\footnotesize
\begin{tabularx}{\textwidth}{|>{\raggedright\arraybackslash}p{2.3cm}|X|}
\hline
\textbf{Tool} & \textbf{Purpose} \\
\hline
cuBLAS & GPU-accelerated linear algebra \\
\hline
cuFFT & FFT library \\
\hline
cuDNN & Deep-learning kernels \\
\hline
cuda-gdb & Kernel debugging (breakpoints/step) \\
\hline
cuda-memcheck & Memory correctness checks \\
\hline
\end{tabularx}
}

\keysent{Block-size tradeoff from your notes:}\par
{\footnotesize
\begin{tabularx}{\textwidth}{|>{\raggedright\arraybackslash}p{2.9cm}|X|}
\hline
\textbf{Choice} & \textbf{Tradeoff} \\
\hline
\textbf{Larger blocks} & Better shared-memory reuse and fewer redundant global loads; can lower occupancy. \\
\hline
\textbf{Smaller blocks} & Easier latency hiding and scheduling flexibility; less intra-block reuse and potentially lower arithmetic intensity. \\
\hline
\end{tabularx}
}

\subsection{Bottleneck Types}
{\footnotesize
\setlength{\tabcolsep}{2.8pt}
\begin{tabularx}{0.98\textwidth}{|>{\raggedright\arraybackslash}p{2.2cm}|X|X|}
\hline
\textbf{Type} & \textbf{Symptoms} & \textbf{Typical fixes} \\
\hline
Compute-bound & High SM activity, memory bandwidth not maxed, kernels dominate runtime. & Improve Tensor Core use, fuse kernels, improve occupancy, reduce divergence. \\
\hline
Memory-bound & Bandwidth near peak, moderate SM activity, many stalls. & Tiling, shared-memory reuse, coalescing, fewer global reads/writes. \\
\hline
IO-bound & GPU idle, DataLoader maxed, disk/network busy. & Dataset caching, offline preprocessing, parallel data loading, faster storage (SSD/NVMe), reduced contention, distributed sharding with node-local caching. \\
\hline
Communication-bound & Placeholder retained. & Placeholder retained. \\
\hline
\end{tabularx}
}

\keysent{Important CUDA terms:}
\begin{itemize}
\item \texttt{gridDim.x}, \texttt{gridDim.y}: grid dimensions (number of blocks)
\item \texttt{blockDim.x}, \texttt{blockDim.y}: threads per block dimensions
\item \texttt{threadIdx.x}, \texttt{threadIdx.y}: thread coordinates inside block
\item \texttt{blockIdx.x}, \texttt{blockIdx.y}: block coordinates inside grid
\end{itemize}

\keysent{2D indexing reminder (matmul/conv-style):}
\[
i = \texttt{blockIdx.y} \cdot \texttt{blockDim.y} + \texttt{threadIdx.y}, \quad
j = \texttt{blockIdx.x} \cdot \texttt{blockDim.x} + \texttt{threadIdx.x}
\]

\begin{greenbox}{Prompt}
GPU utilization is 40\%. How do you debug?
\begin{enumerate}[leftmargin=1.3em]
\item Confirm low utilization with \texttt{watch -n 1 nvidia-smi}.
\item Check detailed counters: \texttt{nvidia-smi dmon -s u}.
\item Check CPU loader bottlenecks: \texttt{htop}.
\item Check process/thread CPU placement: \texttt{ps -o pid,psr,pcpu,comm -p <PID>}.
\item Check data-loader I/O: \texttt{iostat -xm 1}.
\item Check filesystem/network I/O: \texttt{df -h}, \texttt{lfs df -h}, \texttt{sar -n DEV 1}.
\item Check communication dominance:
\begin{itemize}
\item \texttt{export NCCL\_DEBUG=INFO}
\item \texttt{export NCCL\_DEBUG\_SUBSYS=ALL}
\end{itemize}
\item Profile with Nsight Systems:
\begin{itemize}
\item \texttt{nsys profile -o profile\_report python train.py}
\item look for large NCCL blocks and GPU idle gaps
\end{itemize}
\item Check per-GPU batch size: \(\text{per-GPU batch}=\text{global batch}/N\).
\item Check stragglers and thermals: \texttt{nvidia-smi}, \texttt{nvidia-smi -q -d TEMPERATURE}.
\item Check NUMA/topology and PCIe paths: \texttt{nvidia-smi topo -m}.
\item If kernels are too short/small, increase effective work per step.
\end{enumerate}
\end{greenbox}

\vspace{2cm}

% ======================================================
\section{Data Pipeline Scaling}
% ======================================================

\begin{bluebox}{Topics}
\begin{itemize}
\item Sharding datasets
\item DataLoader \texttt{num\_workers}
\item Prefetching
\item Caching
\item Metadata server overload
\end{itemize}
\end{bluebox}

\keysent{Pipeline details from your notes:}
\begin{itemize}
\item \textbf{DistributedSampler:} avoids duplicate shard reads across GPUs, improves balance across ranks.
\item \textbf{DataLoader workers:} each worker is a subprocess loading data in parallel.
\item \textbf{Prefetching:} loads the next batch while the GPU processes the current batch.
\item Important flags from your notes: \texttt{prefetch\_factor}, \texttt{persistent\_workers}, \texttt{pin\_memory}.
\item \texttt{pin\_memory} allocates page-locked host memory to speed CPU-to-GPU transfers.
\item \textbf{Caching:} cache dataset in RAM, cache preprocessed tensors, or use local NVMe instead of network FS.
\end{itemize}

{\footnotesize
\begin{tabularx}{\textwidth}{|>{\raggedright\arraybackslash}p{3cm}|X|X|}
\hline
\textbf{Issue} & \textbf{Symptoms} & \textbf{Fixes from your notes} \\
\hline
DataLoader starvation & Low GPU utilization; CPU and I/O high; profiling shows gaps between batches. & Increase workers, enable \texttt{pin\_memory}, cache locally, increase batch size, use faster storage, use async preprocessing. \\
\hline
Metadata server overload & Random stalls, high IOPS, inconsistent speed, high \texttt{\%util}/wait in \texttt{iostat -xm 1}. & Use tar/record files, prepack dataset, cache locally, increase striping. \\
\hline
Distributed amplification & Multiple processes load data simultaneously, amplifying I/O pressure and metadata contention. & Shard data, reduce small files, cache node-local. \\
\hline
Slow rank stall & Collectives (for example AllReduce) require all ranks; slowest rank sets step time. & Fix stragglers and rebalance pipeline. \\
\hline
\end{tabularx}
}

\begin{greenbox}{Prompt}
Why might training stall randomly in distributed jobs?
\begin{enumerate}[leftmargin=1.3em]
\item DataLoader starvation.
\item Metadata server overload.
\item Network file system congestion.
\item Straggler node.
\item Random garbage collection or memory spikes.
\item Background system activity (CPU, disk, network).
\end{enumerate}
\end{greenbox}

\vspace{2cm}


\subsection{KV Cache for Autoregressive Inference}

\begin{bluebox}{KV Cache Mental Model}
In decoder-only inference, the model generates one token at a time.  
For each new token, you only need a \textbf{new query} and the historical \textbf{cached keys/values} from prior tokens.
\end{bluebox}

From your notes, the key idea is to avoid recomputing old attention state repeatedly:
\begin{itemize}
\item Transformer layers produce \(Q, K, V\) per token.
\item For the newest token, attention uses that token's \(Q_t\) against all previous \(K_{1:t}\) and \(V_{1:t}\).
\item Recomputing \(K,V\) for old tokens every step is redundant.
\end{itemize}

\textbf{KV-cache decode loop (cleaned from your notes):}
\begin{enumerate}[leftmargin=1.3em]
\item Compute \(Q_t, K_t, V_t\) for the newest token only.
\item Append \(K_t, V_t\) to per-layer cache.
\item Read cached \(K_{1:t}, V_{1:t}\).
\item Compute attention with \(Q_t\) and cached history.
\end{enumerate}

\begin{redbox}{Important Clarification}
KV caching removes repeated \textbf{projection} work for past tokens, but attention over context is still sequence-length dependent.  
So decode compute does not become constant-time; it becomes much more efficient than full recomputation.
\end{redbox}

\textbf{TTFT vs steady decode (from your notes):}
\begin{itemize}
\item \textbf{Prefill phase:} process full prompt, build cache, highest initial compute cost.
\item \textbf{Decode phase:} one-token steps using warm cache.
\item Longer prompts increase prefill cost, so \textbf{time-to-first-token (TTFT)} rises.
\end{itemize}

{\footnotesize
\begin{tabularx}{\textwidth}{|>{\raggedright\arraybackslash}p{2.5cm}|X|X|}
\hline
\textbf{Phase} & \textbf{Main work} & \textbf{Typical bottleneck} \\
\hline
Prefill & Full prompt forward pass + KV cache build & Compute + memory bandwidth \\
\hline
Decode & Per-token forward with cached history & Memory bandwidth + cache capacity \\
\hline
\end{tabularx}
}

\textbf{Memory tradeoff (your point, organized):}
\begin{itemize}
\item KV cache trades compute for memory.
\item Cache grows with layer count, context length, and concurrent requests.
\item For long contexts and many concurrent users, KV memory can exceed model weights.
\item This motivates \textbf{MQA/GQA} to reduce KV-head memory pressure.
\end{itemize}



% ======================================================
\section{Fault Tolerance + Checkpointing}
% ======================================================

\keysent{How it is done in PyTorch (from your notes): use DistributedSampler.}

\begin{lstlisting}[language=Python, caption={DistributedSampler sketch from notes}]
sampler = torch.utils.data.DistributedSampler(
    dataset,
    num_replicas=world_size,
    rank=rank,
    shuffle=True,
)
loader = DataLoader(dataset, sampler=sampler)
\end{lstlisting}

\keysent{Why this matters (from your notes):}
\begin{itemize}
\item avoids scenarios where GPUs read the same data shard
\item avoids load imbalance across ranks
\item reduces straggler effects from uneven sample counts
\end{itemize}

\keysent{If dataset size is not divisible by world size:}
\begin{itemize}
\item some ranks may process fewer samples, causing sync delays
\item using \texttt{drop\_last=True} can help keep step balance
\end{itemize}

\keysent{DataLoader \texttt{num\_workers} notes:}
\begin{itemize}
\item each worker is a subprocess loading data in parallel
\item too low: GPU waits for data (low utilization)
\item too high: CPU thrashing, memory pressure, context-switch overhead
\item starting point from your notes: around 4--8 workers per GPU (then tune by CPU/core budget)
\item monitor with \texttt{htop}; increase until GPU utilization stops improving or CPU saturates
\end{itemize}

\begin{bluebox}{Prefetching + Caching Notes}
\begin{itemize}
\item Prefetching loads the next batch while the GPU processes the current batch.
\item Without prefetching, GPU idle gaps appear between batches.
\item Shared filesystem failures from your notes: metadata overload, thousands of small file opens, random stalls.
\item Small files hurt because each access requires metadata lookup and adds latency.
\item Metadata overload (Lustre/NFS): random stalls, high IOPS, inconsistent speed.
\item Detection: \texttt{iostat -xm 1} (high utilization and high wait time).
\item Fixes from your notes: tar/record files, prepack dataset, cache locally, increase striping.
\end{itemize}
\end{bluebox}

\keysent{Possible questions from your notes:}\par
{\footnotesize
\begin{tabularx}{\textwidth}{|>{\raggedright\arraybackslash}p{3.4cm}|X|}
\hline
\textbf{Question} & \textbf{Answer from your notes (cleaned)} \\
\hline
How do you detect DataLoader bottleneck? & GPU util low, CPU usage high, I/O high, profiling shows gaps between batches. \\
\hline
How would you improve data pipeline throughput? & Increase workers, enable pin memory, cache locally, increase batch size, use faster storage, use async preprocessing. \\
\hline
Why does distributed training amplify data bottlenecks? & Multiple processes load data simultaneously, increasing I/O pressure and metadata contention. \\
\hline
Why does one slow rank stall everything? & Collective operations like AllReduce require all ranks; the slowest process determines global step time. \\
\hline
\end{tabularx}
}

\subsection{Checkpoint and Resume State}
\begin{itemize}
\item \textbf{Sharded checkpoints}
\item \textbf{Partial node failure}
\item \textbf{Resume logic}
\item \textbf{Optimizer state restoration}
\item \textbf{RNG state}
\end{itemize}

\keysent{Checkpointing details from your notes:}
\begin{itemize}
\item Checkpointing can cause memory spikes if consolidation gathers full parameters temporarily.
\item Reduce checkpoint time with sharded saves, fewer large files, less frequent saves, and async writing.
\item For every parameter, optimizer stores momentum and variance; optimizer memory can be around \(2\times\) model size.
\item Also save step count and learning-rate scheduler state.
\item Restoring only weights (without optimizer state) loses optimizer history and can behave like partial reset.
\end{itemize}

\keysent{What to save (minimum from your notes):}
\begin{itemize}
\item model weights (sharded)
\item optimizer state (sharded)
\item global step
\item LR scheduler state
\item RNG state
\item ideally sampler state
\end{itemize}

\keysent{RNG APIs in your notes:}
\begin{itemize}
\item \texttt{torch.get\_rng\_state()}
\item \texttt{torch.cuda.get\_rng\_state\_all()}
\item \texttt{numpy.random.get\_state()}
\item \texttt{random.getstate()}
\end{itemize}

\keysent{Why this matters:} dropout and data shuffling depend on RNG, so resume is not deterministic without restoring RNG state.

\begin{redbox}{Hard Question}
How do you checkpoint a 70B model efficiently?

\begin{itemize}
\item do not gather full model to one GPU
\item save sharded checkpoints
\item save optimizer state
\item save training step + RNG state
\item write fewer large files, not thousands of small files
\end{itemize}

\keysent{Scale note from your content:}
\begin{itemize}
\item in BF16, weights alone can be very large (your note: about 140 GB)
\item Adam adds momentum and variance states, so total distributed training state can be much larger (your note: around 400 GB)
\item gathering to one GPU can lead to OOM
\end{itemize}

\keysent{What to implement from your notes:}
\begin{itemize}
\item in FSDP or ZeRO-3, each GPU saves its parameter shard and optimizer shard
\item use sharded checkpointing
\item avoid consolidation unless exporting for inference
\end{itemize}
\end{redbox}

\vspace{2cm}

% ======================================================
\section{Mixed Precision + Numerical Stability}
% ======================================================

\subsection{FP16 vs BF16}
{\footnotesize
\begin{tabularx}{\textwidth}{|>{\raggedright\arraybackslash}p{2.2cm}|X|X|}
\hline
\textbf{Format} & \textbf{Bit layout from your notes} & \textbf{Key point} \\
\hline
FP16 & 16 bits total, 5-bit exponent, 10-bit mantissa & Smaller dynamic range; more prone to underflow/overflow. \\
\hline
BF16 & 16 bits total, 8-bit exponent, 7-bit mantissa & FP32-like exponent range; more stable in practice. \\
\hline
\end{tabularx}
}

\begin{itemize}
\item Why BF16 is more stable: same exponent range as FP32, so large and small numbers are represented more safely.
\item Why FP16 may diverge: gradients can underflow to zero or overflow to infinity due to limited exponent range.
\end{itemize}

\subsection{Loss Scaling}
\begin{itemize}
\item Needed because FP16 has a small dynamic range.
\item Very small gradients may underflow to zero.
\item Scale loss before backward, then divide gradients by the scale after backward.
\item \textbf{Static scaling:} constant scale.
\item \textbf{Dynamic scaling:} preferred in your notes; adjusts scale based on overflow detection.
\end{itemize}

\subsection{Gradient Overflow}
\begin{itemize}
\item Overflow happens when gradient values exceed representable range.
\item Symptoms: loss becomes NaN, gradients contain \texttt{inf}, training diverges suddenly.
\item Detection from your notes:
\begin{itemize}
\item \texttt{torch.autograd.detect\_anomaly()}
\item GradScaler checks
\item \texttt{torch.isnan(tensor).any()}
\end{itemize}
\item Causes from your notes: large activations, unstable learning rates, insufficient loss scaling in FP16.
\item Debug from your notes: check overflow, reduce learning rate, enable gradient clipping, switch to BF16, inspect layer outputs.
\end{itemize}

\subsection{Stability Tradeoffs}
\begin{itemize}
\item \textbf{Memory savings:} 16-bit instead of 32-bit gives roughly \(2\times\) memory savings.
\item \textbf{Throughput increase:} Tensor Cores are optimized for FP16/BF16.
\item \textbf{Numerical risk:} reduced precision can increase rounding and accumulation error.
\end{itemize}

\keysent{Typical strategy from your notes:}
\begin{itemize}
\item store master weights in FP32
\item compute in FP16 or BF16
\item reduce in FP16 or BF16
\item update in FP32
\end{itemize}

\keysent{Common Q\&A from your notes:}
\begin{itemize}
\item Why keep master weights in FP32? To preserve update precision and reduce accumulated rounding error.
\item When avoid mixed precision? Unstable models, small-model debugging, or unsupported hardware.
\item Why can large batching increase instability risk? Larger batches can increase gradient magnitude and overflow risk at low precision.
\end{itemize}

\vspace{2cm}

% ======================================================
\section{Scheduler + Orchestration}
% ======================================================

\subsection{Slurm}
\textbf{Slurm} stands for \textit{Simple Linux Utility for Resource Management}.

\begin{itemize}
\item \texttt{sbatch} vs \texttt{srun}
\begin{itemize}
\item \texttt{sbatch}: non-interactive, queued, runs when resources are available; commonly used for long training jobs, production runs, and multi-node jobs.
\item \texttt{srun}: launches tasks; interactive or used inside allocation; common for launching distributed training processes.
\end{itemize}

\item Partitions
\begin{itemize}
\item Logical queues (for example: gpu, cpu, highmem, debug).
\item Configured in \texttt{slurm.conf}.
\item Each partition has time limits, node restrictions, and resource limits.
\item Used to separate workloads, enforce policies, and control access for different job types.
\end{itemize}

\item GPU allocation
\begin{itemize}
\item In Slurm:
\texttt{\#SBATCH --gres=gpu:4}
\item GRES = generic resources.
\item Slurm assigns GPU IDs and sets \texttt{CUDA\_VISIBLE\_DEVICES}.
\item Multi-node example:
\texttt{\#SBATCH --nodes=2}\\
\texttt{\#SBATCH --ntasks-per-node=4}\\
\texttt{\#SBATCH --gres=gpu:4}
\end{itemize}

\item Fairshare
\begin{itemize}
\item Slurm tracks historical usage.
\item If you used many GPUs recently, your priority can drop.
\item Fairshare improves fairness and prevents monopolization.
\item Priority factors: fairshare, job age, QoS, partition priority.
\end{itemize}
\end{itemize}

\subsection{Kubernetes}
Here scheduling is for containers, not raw processes.
Typical use from your notes: cloud GPU clusters, microservices, ML platforms (EKS/GKE/AKS).

\begin{itemize}
\item Pods
\begin{itemize}
\item Smallest deployable unit.
\item Contains one or more containers with shared network and storage.
\end{itemize}
\begin{lstlisting}[language=yaml, caption={Pod GPU limit sketch from notes}]
resources:
  limits:
    nvidia.com/gpu: 1
\end{lstlisting}

\item GPU resource requests
\begin{itemize}
\item Kubernetes uses a device plugin for GPU resources.
\end{itemize}
\begin{lstlisting}[language=yaml, caption={Kubernetes GPU request/limit sketch from notes}]
resources:
  requests:
    nvidia.com/gpu: 2
  limits:
    nvidia.com/gpu: 2
\end{lstlisting}
\begin{itemize}
\item Scheduler assigns a node with enough GPUs and exposes them to the container.
\item Kubernetes knows GPU resources via device-plugin registration.
\end{itemize}

\item Scheduling policies
\begin{itemize}
\item Scheduler considers CPU, memory, GPU, node labels, taints/tolerations, and affinity.
\end{itemize}
\begin{lstlisting}[language=yaml, caption={Node selector sketch from notes}]
nodeSelector:
  gpu-type: a100
\end{lstlisting}

\item Auto-scaling
\begin{itemize}
\item Kubernetes can scale pods and clusters.
\item Useful for inference and dynamic training environments.
\end{itemize}
\end{itemize}

\keysent{Q\&A from your notes:}\par
{\footnotesize
\begin{tabularx}{\textwidth}{|>{\raggedright\arraybackslash}p{4.1cm}|X|}
\hline
\textbf{Question} & \textbf{Answer from your notes (cleaned)} \\
\hline
Why is autoscaling useful for ML? & Inference demand fluctuates, and training clusters may need to add/remove GPU nodes dynamically. \\
\hline
Slurm vs Kubernetes for ML workloads? & Slurm: HPC-oriented, better for tightly coupled multi-node training, native MPI/NCCL integration, lower overhead. Kubernetes: container-native, easier cloud integration, built-in autoscaling, flexible workload types. \\
\hline
Why might distributed training hang under Slurm? & Wrong environment variables, firewall between nodes, incorrect tasks-per-node, misconfigured NCCL. \\
\hline
How does Slurm integrate with DDP? & Slurm sets env vars like \texttt{SLURM\_PROCID} and \texttt{SLURM\_NODEID} that can initialize rank and world size. \\
\hline
Why might Kubernetes be worse for tightly coupled training? & Higher networking overhead, container abstraction layers, and less optimized HPC fabric integration compared to Slurm. \\
\hline
How would you schedule multi-node DDP on Kubernetes? & Use StatefulSets or custom job controllers with proper rank assignment and headless services for communication. \\
\hline
What is gang scheduling? & Allocate all resources for a distributed job simultaneously to avoid partial scheduling. \\
\hline
\end{tabularx}
}

\vspace{2cm}

% ======================================================
\section{Storage Systems}
% ======================================================

\begin{bluebox}{Storage Layers from your notes}
\begin{itemize}
\item local NVMe (fast per node)
\item parallel filesystem (Lustre/GPFS/BeeGFS)
\item possible object storage (S3-like)
\end{itemize}
For large-scale ML, shared parallel storage is critical.
\end{bluebox}

\begin{itemize}
\item Lustre architecture
Lustre is a parallel distributed filesystem used in HPC.

It separates metadata and data storage.

\textbf{Components:}
\begin{enumerate}[leftmargin=1.4em]
\item \textbf{MDS (metadata server)}: stores filenames, permissions, directory structure, and file-location mapping; handles \texttt{open()}, \texttt{close()}, \texttt{stat()} metadata operations.
\item \textbf{MDT (metadata target)}: storage device behind MDS.
\item \textbf{OSS (object storage server)}: handles actual file data.
\item \textbf{OST (object storage target)}: stores file data blocks.
\end{enumerate}

Files are striped across multiple OSTs for parallel access.

What happens when you open a file?
\begin{itemize}
\item client asks MDS for metadata
\item MDS returns file layout
\item client reads data directly from OST
\end{itemize}

Q: Why separate metadata and data?  
To allow parallel data access across multiple storage targets while centralizing lighter metadata operations.

Q: What happens when metadata server is overloaded?  
File open/stat operations slow down, causing random stalls even when data bandwidth is fine.

\item Metadata server
This is where many ML workloads fail.

ML datasets often contain:
\begin{itemize}
\item millions of small image files
\item many file opens per batch
\item multiple DataLoader workers per GPU
\end{itemize}

Each worker may call \texttt{open()} thousands of times. If many GPUs do this simultaneously, metadata storms occur.

Symptoms:
\begin{itemize}
\item random training stalls
\item high IOPS but low throughput
\item GPU utilization fluctuates
\item CPU waits on I/O
\end{itemize}

Detection:
\begin{itemize}
\item \texttt{iostat -xm 1} (look for high \texttt{\%util} and high wait)
\item \texttt{lfs df -h} on Lustre
\end{itemize}

Q: Why do small files hurt parallel filesystems?  
Each file access needs metadata lookup, which can overwhelm the metadata server.

Q: Why might training stall randomly?  
Metadata server bottlenecks under high concurrent file-open rates.

\item Small file problem
Very common in ML.

Example from your notes: ImageNet has many small JPEG files.

Per batch, many files are opened repeatedly across many GPUs.

Problem:
\begin{itemize}
\item high metadata overhead
\item poor sequential read efficiency
\item high open/close syscall overhead
\end{itemize}

Solutions from your notes:
\begin{itemize}
\item pack into larger containers: TFRecord, WebDataset, LMDB, HDF5
\item increase striping: \texttt{lfs setstripe -c 4 myfile}
\item cache locally (for example: \texttt{rsync -av dataset /local\_nvme/})
\end{itemize}

Q: How to fix the small-file problem?  
Batch files into larger shard files to reduce metadata lookups and improve sequential I/O.

\item Parallel IO
Lustre allows files to be striped across multiple OSTs so clients can read in parallel.

Striping improves bandwidth; \texttt{lfs getstripe filename} shows file striping.

Why parallel I/O matters:
\begin{itemize}
\item large model checkpointing (100+ GB)
\item many GPUs reading simultaneously
\item need aggregate bandwidth
\end{itemize}

Without striping:
\begin{itemize}
\item single-OST bottleneck
\item slower checkpoint save/load
\end{itemize}

What is file striping?  
Dividing file data across multiple storage targets to enable parallel read/write.

When would you increase stripe count?  
For large files like checkpoints to increase aggregate throughput.

Common GPU cluster storage failures from your notes:
\begin{itemize}
\item GPUs at 40\% utilization caused by data pipeline I/O bottlenecks and metadata overload
\item checkpoints taking minutes due to single-OST bottleneck, insufficient striping, or network congestion
\item random stalls due to shared filesystem contention and too many simultaneous jobs
\end{itemize}

Q: Why might local NVMe be faster than Lustre?  
It removes network overhead and metadata contention, giving lower latency and more consistent per-node throughput.

Q: Why is a parallel filesystem needed?  
Multi-node training needs shared access to datasets and checkpoints.

Q: Why can distributed training amplify I/O issues?  
Many ranks read simultaneously, multiplying metadata and bandwidth pressure.

Q: Why does checkpointing stall other jobs?  
Large parallel writes can saturate shared OST bandwidth and metadata operations.

Storage and GPU scaling connection: scaling GPU count increases data-read parallelism, checkpoint write pressure, and metadata operations, so storage becomes a bottleneck at scale.

\end{itemize}

% ======================================================
\section{Cloud vs On-Prem Tradeoffs}
% ======================================================

\begin{itemize}
\item Spot instances
Cloud providers (for example AWS/GCP) offer discounted compute capacity with interruption risk.

Your notes: often around 50--90\% cheaper; useful for hyperparameter sweeps, non-critical jobs, and fault-tolerant workloads.

Risks: instance termination, loss of local storage, and distributed job collapse.

Use spot safely (from your notes): frequent checkpointing, TorchElastic, distributed job managers, persistent object storage for checkpoints.

Q: How would you train a 70B model on spot instances?  
Only if the framework supports fault tolerance and elastic recovery; otherwise interruption risk may waste compute.

Q: How make distributed training resilient to spot termination?  
Use elastic training, periodic checkpointing to persistent storage, and automatic rank rejoin/restart.

\item Cost optimization
On-prem vs cloud cost model from your notes includes GPU hours, storage, bandwidth, data egress fees, idle resources, and networking.

Q: Why might cloud training be more expensive?  
High GPU-hour rates plus large-scale distributed networking/storage costs.

Q: When is on-prem cheaper?  
When GPU utilization is consistently high over long periods and hardware cost is amortized.

Q: Hidden cloud costs in ML?  
Data egress fees and underutilized reserved capacity.

\item Elastic scaling
Ability to add/remove nodes dynamically and adjust cluster size with demand.

Cloud strengths: autoscaling groups, Kubernetes cluster autoscaler, dynamic node provisioning.

On-prem limitation: fixed node pool; cannot instantly add GPUs; needs hardware procurement.

TorchElastic points from your notes:
\begin{itemize}
\item allows ranks to leave/join
\item reforms process groups
\item handles failures
\end{itemize}

Scenario from your notes: training can lose 1--2 workers, restart, and continue from latest checkpoint.

Instead of fixed world size, TorchElastic uses rendezvous coordination, dynamic membership (min/max workers), and automatic restart on failure.

When a worker dies: collectives fail, worker group restarts, ranks/world size reform, and script resumes from checkpoint.

Q: Why is elastic scaling hard for distributed training?  
Most frameworks assume fixed world size and static ranks; changing size requires communication-group changes.

Q: Why is elasticity easier for inference than training?  
Inference is stateless per request, while training requires synchronized gradient updates.

\begin{lstlisting}[language=Python, caption={TorchElastic resume sketch from notes}]
import os
import time
import torch
import torch.distributed as dist
from torch.nn.parallel import DistributedDataParallel as DDP

CHECKPOINT_PATH = "ckpt.pt"

def setup_distributed():
    """
    TorchElastic (torchrun) sets these env vars:
      - RANK, LOCAL_RANK, WORLD_SIZE
      - MASTER_ADDR, MASTER_PORT (or rendezvous handles it)
    """
    dist.init_process_group(backend="nccl")
    local_rank = int(os.environ["LOCAL_RANK"])
    torch.cuda.set_device(local_rank)
    return local_rank

def build_model():
    model = torch.nn.Linear(1024, 1024).cuda()
    return model

def save_checkpoint(ddp_model, optimizer, step: int):
    if dist.get_rank() == 0:
        ckpt = {
            "step": step,
            "model": ddp_model.module.state_dict(),
            "optim": optimizer.state_dict(),
        }
        torch.save(ckpt, CHECKPOINT_PATH)

def load_checkpoint(ddp_model, optimizer):
    if os.path.exists(CHECKPOINT_PATH):
        ckpt = torch.load(CHECKPOINT_PATH, map_location="cpu")
        ddp_model.module.load_state_dict(ckpt["model"])
        optimizer.load_state_dict(ckpt["optim"])
        return int(ckpt["step"])
    return 0

def main():
    local_rank = setup_distributed()
    model = build_model()
    ddp_model = DDP(model, device_ids=[local_rank])
    optimizer = torch.optim.AdamW(ddp_model.parameters(), lr=1e-3)
    start_step = load_checkpoint(ddp_model, optimizer)

    for step in range(start_step, start_step + 10_000):
        optimizer.zero_grad(set_to_none=True)
        x = torch.randn(32, 1024, device="cuda")
        y = ddp_model(x).sum()
        y.backward()
        optimizer.step()
        if step % 200 == 0:
            save_checkpoint(ddp_model, optimizer, step)
        if dist.get_rank() == 0 and step % 200 == 0:
            print(f"[step={step}] world_size={dist.get_world_size()} ok")

    dist.destroy_process_group()

if __name__ == "__main__":
    main()
\end{lstlisting}

Here is the Slurm script to run this:
\begin{lstlisting}[language=bash, caption={Slurm + torchrun launch sketch from notes}]
# inside an allocation with 2 nodes, 4 GPUs/node = 8 procs total
export RDZV_ENDPOINT="$(scontrol show hostnames $SLURM_NODELIST | head -n1):29400"
export NNODES=$SLURM_NNODES

srun bash -c '
  torchrun \
    --nnodes='"$NNODES"' \
    --nproc_per_node=4 \
    --rdzv_backend=c10d \
    --rdzv_endpoint='"$RDZV_ENDPOINT"' \
    --rdzv_id=job123 \
    --max_restarts=10 \
    train_elastic.py
'
\end{lstlisting}

RDZV endpoint is where workers meet, \texttt{max\_restarts} controls tolerated failures, and \texttt{rdzv\_id} is the rendezvous room name that must match.

Q: If a worker dies, does training continue instantly?  
Not usually. Current process group breaks; TorchElastic restarts workers and reforms ranks; resume happens via checkpoint load.

Q: What do you need to checkpoint for correct resume?  
Minimum: model weights, optimizer state, global step. Ideally: RNG state and sampler state.

Q: Does world size change automatically?  
It can, depending on min/max config; many setups keep world size constant and only restart on failures.

Biggest pitfall from your notes: script restarts from scratch unless resume-from-checkpoint logic is implemented.

With autoscaling, a job can start at minimum nodes and grow toward max nodes as more capacity arrives.

\item Network topology differences
Uses AWS Elastic Fabric Adapter (EFA) instead of InfiniBand in many cloud setups.

Data is often in S3, but large multi-node training is usually better when data is staged locally or served by a high-performance filesystem to avoid high latency.

Uses Scalable Reliable Datagram (SRD), designed to handle congestion via multipath networking.

Large-scale training depends on AllReduce bandwidth, low latency, and network determinism; cloud can introduce higher jitter and bandwidth variability.

Jitter is variability in network latency. It causes synchronization delays and amplifies straggler effects in distributed training.

\end{itemize}

% ======================================================
\section{Researcher Productivity Metrics}
% ======================================================

{\footnotesize
\begin{tabularx}{\textwidth}{|>{\raggedright\arraybackslash}p{2.8cm}|X|X|}
\hline
\textbf{Metric} & \textbf{What to measure from your notes} & \textbf{Interpretation from your notes} \\
\hline
GPU utilization & \texttt{nvidia-smi}, \texttt{nvidia-smi dmon}, \texttt{nsys}. & Under ~40\% often underperforming; around ~90\% can indicate healthy compute-bound behavior (but still check queue/convergence). \\
\hline
Queue time & submission \(\rightarrow\) job start; track average and P95 wait time. & Lower queue time improves researcher iteration speed. \\
\hline
Time-to-convergence & wall-clock time to target accuracy/loss (not epochs). & Most important end metric; throughput alone is insufficient. \\
\hline
Experiment reproducibility & controlled RNG, stable environment, deterministic ops, proper checkpointing. & Without reproducibility, debugging and conclusions are unreliable. \\
\hline
Cost per training run & \(\text{cost}=\text{GPU hours}\times\text{cost per GPU hour}\) + storage/network/engineering. & Key for cost per experiment/model/inference impact. \\
\hline
\end{tabularx}
}

\keysent{Notes from your content:}
\begin{itemize}
\item Low utilization causes: DataLoader bottleneck, network bottleneck, small batch size, poor overlap, CPU starvation.
\item Reduce queue time with scheduling policy improvements, fairshare tuning, dedicated partitions, preemption, and cloud elasticity.
\item If throughput improves but time-to-convergence does not: check scaling efficiency, communication bottlenecks, LR scaling, and batch-size optimization effects.
\item Improve reproducibility: fix seeds, log environment, version data/code, save full training state.
\item Reduce run cost: increase scaling efficiency, use mixed precision, optimize batch size, reduce failed runs, use spot safely.
\end{itemize}

\begin{greenbox}{Prompt}
How do you measure infrastructure impact?
\begin{enumerate}[leftmargin=1.3em]
\item Define baseline:
\begin{itemize}
\item GPU utilization
\item queue time
\item step time
\item time to convergence
\item failure rate
\item cost per run
\end{itemize}
\item Apply infrastructure improvements:
\begin{itemize}
\item better networking
\item improved data loader
\item FSDP adoption
\item storage caching
\end{itemize}
\item Remeasure and compare:
\begin{itemize}
\item percent improvement in step time
\item reduction in queue time
\item reduction in job failures
\item reduction in cost per convergence
\end{itemize}
\end{enumerate}
\end{greenbox}

\vspace{2cm}
\section{Rapid-Fire Questions Section}
% ======================================================

\begin{bluebox}{Rapid-Fire: One-Line Answer Bank}
\begin{itemize}
\item \textbf{Why doesn't scaling linearly improve speed?} Distributed training adds communication and synchronization overhead (for example, AllReduce). As GPU count grows, network contention, latency amplification, and straggler effects grow too. If global batch size stays constant, per-GPU compute shrinks and communication-to-compute ratio worsens, so scaling becomes network-bound.

\item \textbf{What is gradient accumulation?} Gradients from multiple forward/backward micro-steps are accumulated before one optimizer step. This simulates a larger effective batch without requiring a larger per-step memory footprint.

\item \textbf{What is microbatching?} A large batch is split into smaller microbatches processed sequentially (often in pipeline parallelism). This keeps devices busier and reduces pipeline idle time, but increases scheduling complexity.

\item \textbf{What causes NCCL deadlocks?} Mismatched collective calls across ranks, crashed ranks, wrong world-size/rank initialization, firewall/network misconfiguration, or version/runtime mismatch. One misbehaving rank can stall the job.

\item \textbf{Why do small files hurt Lustre?} Each file open requires metadata operations. With many workers and many GPUs, metadata requests explode and overload the metadata server, causing random stalls and low GPU utilization.

\item \textbf{Why does large batch size hurt generalization?} Large batches reduce gradient noise and can converge to sharper minima that generalize worse. They also require correct learning-rate scaling.

\item \textbf{What is straggler effect?} One slow worker delays all others at synchronization points, so overall step time is determined by the slowest rank.

\item \textbf{What is NUMA?} Non-uniform memory access: memory latency depends on CPU socket locality. Remote-memory/PCIe access paths can hurt performance.

\item \textbf{What is PCIe bottleneck?} Host-device or inter-device transfers saturate PCIe bandwidth, limiting throughput when communication/memory movement dominates.

\item \textbf{What is kernel fusion?} Multiple operations are combined into one kernel launch to reduce memory traffic and launch overhead, improving locality and throughput.

\item \textbf{What is compute intensity?} Arithmetic operations per byte transferred (FLOPs/byte). High intensity is usually compute-bound; low intensity is usually memory-bound.

\item \textbf{Why does communication dominate at scale?} As rank count rises, collective overhead grows while per-rank compute can shrink, so synchronization time can dominate step time.

\item \textbf{What happens if one rank crashes?} In synchronous training, collectives fail/hang and the job typically stops. In elastic setups, the worker group restarts and resumes from checkpoint.

\item \textbf{What is asynchronous training?} Workers update parameters without strict step-by-step synchronization. Throughput can increase, but gradient staleness may hurt convergence stability.

\item \textbf{What is pipeline bubble?} Idle time at pipeline fill/drain phases in pipeline parallel training. Bubble overhead increases when stage count is high relative to microbatch count.

\item \textbf{What is memory fragmentation?} Repeated allocations/deallocations leave non-contiguous free memory blocks; OOM can occur despite enough total free memory.

\item \textbf{What is KV cache?} During autoregressive inference, key/value tensors from previous tokens are stored and reused to avoid recomputing full attention each step.

\item \textbf{What is tensor parallel shard?} A partition of large tensors/weights across GPUs (row-wise or column-wise). Each GPU computes partial outputs, then results are gathered/reduced.

\item \textbf{What is gradient clipping?} Limits gradient norm/magnitude to prevent exploding gradients and improve training stability.

\item \textbf{What is elastic training?} Training that tolerates worker failures/membership changes via rendezvous and restart logic, typically resuming from checkpoints.
\end{itemize}
\end{bluebox}

\vspace{2cm}

% ======================================================
\section{Whiteboard System Design Prompts}
% ======================================================

\begin{greenbox}{Design Prompts}
\begin{itemize}
\item Design multi-tenant GPU scheduling system
\item Design distributed LLM inference cluster
\item Design hyperparameter search system
\item Design fault-tolerant distributed training
\item Design GPU utilization monitoring dashboard
\end{itemize}
\end{greenbox}

\subsection{Design multi-tenant GPU scheduling system}
\keysent{Clarify first:} on-prem vs cloud vs hybrid.

\textbf{Core design from your notes:}
\begin{enumerate}[leftmargin=1.3em]
\item Fairness and quotas: guaranteed quota + burst quota per team; weighted fairshare/DRF-style policy to prevent monopolization.
\item Scheduling strategy: two-level scheduling (choose tenant, then highest-priority job in tenant). Support interactive, production inference, and batch classes.
\item Placement strategy: topology-aware placement; prefer same NVLink domain/same node before cross-node. Use gang scheduling for multi-node distributed jobs.
\item Preemption and isolation: preempt best-effort jobs with checkpoint grace periods; enforce GPU/network/storage isolation per tenant.
\item Observability: expose clear pending reasons (quota exceeded, waiting for gang allocation, topology constraints).
\end{enumerate}

\textbf{Questions to expect:}
\begin{itemize}
\item How do you prevent fragmentation? Bin-pack small jobs and reserve capacity pools for 8--16 GPU gang allocations.
\item How do you handle distributed jobs? Gang scheduling so all required GPUs are allocated atomically.
\item How do you keep utilization high? Backfilling + burst capacity when quotas are underused.
\item Hardest problem? Balancing fairness and utilization while avoiding fragmentation and starvation.
\end{itemize}

\subsection{Design distributed LLM inference cluster}
\keysent{Clarify first:} latency target (P50/P99), tokens/sec, traffic shape, model sizes, hardware, and multi-tenancy constraints.

\textbf{Architecture from your notes:}
\begin{itemize}
\item LLM-aware router in front of GPU worker groups.
\item Continuous batching + KV cache management.
\item Admission control by tenant quotas and estimated KV-memory footprint.
\item Route by cache locality and headroom.
\item Use tensor parallelism for larger models; combine TP within node + pipeline across nodes when needed.
\item Reliability via health checks, circuit breakers, graceful draining, and canary rollouts.
\item Observability: TTFT, per-token latency, tokens/sec, KV eviction/swap rates.
\end{itemize}

\textbf{Questions to expect:}
\begin{itemize}
\item Main bottleneck? Often memory bandwidth + KV cache capacity/management.
\item How keep P95 low with batching? Separate interactive pools, cap token budgets, prioritize short requests.
\item What if KV cache does not fit? Paged KV cache and controlled CPU/NVMe offload fallback.
\item How do you scale? Add model-group replicas and autoscale per model and priority tier.
\end{itemize}

\subsection{Design hyperparameter search system}
\keysent{Clarify first:} objective (best quality, fastest time-to-best, or cost cap), workload type, search space, cluster constraints, and reporting metrics.

\textbf{Core architecture from your notes:}
\begin{itemize}
\item Orchestrator proposes trials.
\item Scheduler launches trials (Slurm/Kubernetes) with quotas and fairshare.
\item Workers report metrics and checkpoints.
\item Default algorithm: ASHA/Hyperband for scalable early stopping; optional TPE for expensive trials.
\item Intermediate metrics stream back for early stop and reprioritization.
\item Checkpoint/resume is first-class for preemption/spot.
\item Centralized stores for trial configs, metrics, artifacts, reproducibility metadata, and UI ranking.
\end{itemize}

\textbf{Questions to expect:}
\begin{itemize}
\item Why random search over grid? In high dimensions, random exploration is often more efficient.
\item Why ASHA? Massively parallel and saves GPU time by early-stopping bad trials.
\item How handle noisy metrics? Smoothing, minimum-step rules, repeated runs for top candidates.
\item Biggest bottleneck? GPU availability/scheduling and storage metadata pressure if checkpointing too frequently.
\item How avoid wasting GPUs? Early stopping, OOM/NaN detection, health checks, automatic retry/resume.
\end{itemize}

\subsection{Design fault-tolerant distributed training}
\textbf{Design summary from your notes:}
\begin{itemize}
\item Treat training as a restartable worker group.
\item Use consistent sharded checkpoints.
\item Add watchdog/heartbeat-based failure and hang detection.
\item Use elastic rendezvous so workers can reform and resume from latest committed checkpoint.
\item Save optimizer state + RNG state; avoid full-state gather that can cause OOM.
\end{itemize}

\textbf{Questions to expect:}
\begin{itemize}
\item What if one rank crashes in DDP? Collectives fail/hang; job restarts in elastic setup.
\item How checkpoint FSDP/ZeRO-3 efficiently? Save per-rank shards; avoid full gather; include optimizer/scheduler/RNG.
\item How avoid corrupted checkpoints? Two-phase commit (temp + manifest + atomic done marker).
\item How detect hangs vs slow steps? Heartbeat/progress thresholds + correlate with NCCL logs and GPU utilization.
\item Biggest pitfall? Not saving optimizer/RNG state; full checkpoint consolidation causing OOM; misdiagnosing DataLoader vs NCCL bottlenecks.
\end{itemize}

\subsection{Design GPU utilization monitoring dashboard}
\textbf{Goal from your notes:}
Build a Grafana dashboard backed by Prometheus, enriched with Slurm/Kubernetes metadata and application training metrics.

\textbf{Core dashboard views:}
\begin{itemize}
\item SM utilization
\item NCCL time
\item DataLoader I/O time
\item throttling/thermal signals
\item tenant-level waste attribution
\end{itemize}

\textbf{Questions to expect:}
\begin{itemize}
\item How map GPU utilization to a specific job? Join scheduler allocation metadata (job \(\rightarrow\) node \(\rightarrow\) GPU IDs) with per-GPU metrics labels.
\item Most actionable metrics? SM utilization + NCCL time + DataLoader time.
\item What alerts? GPUs idle while allocated, ECC/Xid spikes, thermal throttling, NCCL failures, filesystem latency spikes, sustained low utilization by tenant.
\item How avoid metric-overload cost? 1--5s scrape for critical GPU metrics, downsample long-term, and selectively enable high-cardinality labels.
\end{itemize}

\subsection{Simulation 1: Multi-Tenant Heavy Transformer Training (Communication/Network Bottleneck)}

\begin{greenbox}{Scenario from your notes}
Scale from 8 GPUs on one NVLink node (\(\sim85\%\) efficiency) to 32 GPUs across 4 InfiniBand-connected nodes, and efficiency drops to \(\sim60\%\) with high GPU idle time.
\end{greenbox}

\textbf{Diagnostic flow (cleaned):}
\begin{enumerate}[leftmargin=1.3em]
\item Check coarse utilization first (\texttt{nvidia-smi}, DCGM): low SM util + high wait usually means comm or input bottleneck.
\item Use deterministic profiling (\texttt{nsys} + PyTorch profiler): quantify NCCL time vs compute time.
\item Inspect NCCL behavior:
\begin{itemize}
\item \texttt{NCCL\_DEBUG=INFO}, \texttt{NCCL\_DEBUG\_SUBSYS=ALL}
\item verify topology and fabric path (\texttt{nvidia-smi topo -m})
\item confirm traffic is actually on InfiniBand/RDMA, not fallback paths
\end{itemize}
\item Check batch math:
\[
\text{per-GPU batch}=\frac{\text{global batch}}{\#\text{GPUs}}
\]
If global batch is fixed while GPU count increases, per-GPU compute shrinks and communication dominates.
\end{enumerate}

\begin{redbox}{Why Scaling Stops Being Linear}
Synchronous training has barrier-like behavior at collectives.  
One slower rank (CPU jitter, loader delay, network tail latency) can stall all others, so end-to-end step time tracks the slowest rank.
\end{redbox}

\textbf{Mitigations from your notes (organized):}
\begin{itemize}
\item Increase global batch (or use gradient accumulation) to improve compute/communication ratio.
\item Tune DDP bucket size for overlap of backward compute and communication.
\item Use hierarchical collectives for multi-node jobs (intra-node first, then inter-node).
\item Keep TP-heavy communication mostly intra-node when possible.
\end{itemize}

{\footnotesize
\begin{tabularx}{\textwidth}{|>{\raggedright\arraybackslash}p{2.8cm}|X|X|}
\hline
\textbf{Transport} & \textbf{Path cost} & \textbf{Training impact} \\
\hline
TCP/IP sockets & More host copies + kernel/network stack overhead & Higher latency and CPU overhead at scale \\
\hline
RDMA (InfiniBand/EFA) & NIC offload + lower software overhead data path & Better collective latency/bandwidth behavior \\
\hline
\end{tabularx}
}

\subsection{Simulation 2: Memory Monster (100B model on 64 A100 GPUs)}

\begin{greenbox}{Scenario from your notes}
Design distributed training for a 100B-parameter model. The main problem is memory capacity, not raw compute throughput.
\end{greenbox}

\textbf{Why vanilla DDP fails (cleaned math from your notes):}
\begin{itemize}
\item FP16 weights: \(2\) bytes/param \(\Rightarrow 200\) GB.
\item FP16 gradients: \(2\) bytes/param \(\Rightarrow 200\) GB.
\item FP32 master weights: \(4\) bytes/param \(\Rightarrow 400\) GB.
\item Adam states (\(m,v\)): \(8\) bytes/param \(\Rightarrow 800\) GB.
\end{itemize}
Total model-state footprint \(\approx 16\) bytes/param \(\Rightarrow\) \(\sim1.6\) TB (excluding activation peaks).

\begin{redbox}{Implication}
No single 80 GB A100 can hold this full training state.  
You need sharding and model-parallel strategies, not plain replicated DDP.
\end{redbox}

\textbf{Strategy stack from your notes:}
\begin{itemize}
\item \textbf{TP:} shard large matrix multiplies (usually best intra-node on NVLink/NVSwitch).
\item \textbf{PP:} split layers into stages; manage bubble overhead with microbatching.
\item \textbf{ZeRO/FSDP:} shard optimizer, gradients, and parameters to reduce per-GPU steady-state memory.
\end{itemize}

\textbf{ZeRO stage intuition (as you described):}
\begin{itemize}
\item ZeRO-1: shard optimizer states.
\item ZeRO-2: shard optimizer states + gradients.
\item ZeRO-3/FSDP FULL\_SHARD: shard optimizer states + gradients + parameters.
\end{itemize}

\[
\text{Model-state memory per GPU (rough)} \;\propto\; \frac{1}{N}
\]
for fully sharded states across \(N\) ranks, with increased communication.

\textbf{Checkpoint failure mode from your notes:}
\begin{itemize}
\item Training is stable, but periodic checkpoint save causes OOM.
\item Root cause: accidental consolidation/all-gather of full model state onto one rank/device.
\end{itemize}

\textbf{Fix (organized):}
\begin{itemize}
\item Use distributed/sharded state dict checkpointing.
\item Each rank writes its own parameter and optimizer shards.
\item Save full resume context: model shards, optimizer shards, global step, LR scheduler, RNG state.
\end{itemize}

\subsection{Simulation 3: Data Starvation + Metadata Storm}

\begin{greenbox}{Scenario from your notes}
Train a large vision model with hundreds of millions of small JPEG files.  
Observed: SM utilization around \(40\%\), CPU near \(100\%\), high I/O wait, low effective storage throughput.
\end{greenbox}

\textbf{Diagnosis (organized):}
\begin{itemize}
\item Low InfiniBand utilization and high CPU/I/O wait suggest input pipeline bottleneck, not NCCL.
\item Too many small-file \texttt{open/stat} operations overwhelm metadata services.
\item DataLoader workers stall waiting for metadata/file opens, then GPU queues drain.
\end{itemize}

\textbf{Why this happens on Lustre-like systems (from your notes):}
\begin{itemize}
\item Metadata (MDS/MDT) and data (OSS/OST) are separate paths.
\item Small files create heavy metadata traffic before data reads can begin.
\item High fan-out from many nodes \(\times\) many workers can create a metadata storm.
\end{itemize}

\begin{redbox}{Important Clarification}
\texttt{pin\_memory=True} improves host-to-device transfer efficiency, but it does \textbf{not} solve metadata bottlenecks caused by millions of small files.
\end{redbox}

\textbf{Triage steps from your notes:}
\begin{enumerate}[leftmargin=1.3em]
\item Tune \texttt{num\_workers} empirically (avoid over-saturating CPU and metadata path).
\item Enable \texttt{pin\_memory=True} for faster/asynchronous H2D staging.
\item Use \texttt{prefetch\_factor} and \texttt{persistent\_workers=True}.
\item Repack small files into larger sequential shards (WebDataset/TFRecord/LMDB).
\item Add node-local caching on NVMe for repeatedly used datasets.
\end{enumerate}

{\footnotesize
\begin{tabularx}{\textwidth}{|>{\raggedright\arraybackslash}p{3.1cm}|X|X|}
\hline
\textbf{Lever} & \textbf{What it fixes} & \textbf{What it does not fix} \\
\hline
\texttt{pin\_memory=True} & H2D transfer overhead and overlap & Metadata-server overload \\
\hline
Increase \texttt{num\_workers} (carefully) & CPU-side decoding/data prep parallelism & Small-file metadata storm at high worker counts \\
\hline
Shard dataset into large archives & Metadata lookup amplification + random small-file I/O & Poor model/network config issues \\
\hline
Node-local NVMe cache & Network/storage-array contention for reused data & First-load preprocessing cost \\
\hline
\end{tabularx}
}

\textbf{Orchestration choice after pipeline fixes (from your notes):}
\begin{itemize}
\item \textbf{Slurm strengths:} topology-aware placement for tightly coupled synchronous training.
\item \textbf{Kubernetes strengths:} elasticity and cloud-native operations.
\item \textbf{Kubernetes caveat for large distributed training:} usually needs gang scheduling and HPC-network-aware extensions (device plugins, multi-NIC/CNI tuning, scheduler extensions).
\end{itemize}






\vspace{2cm}

% ======================================================
\section{Your Personal Deep Dive Stories}
% ======================================================

\subsection{WaveHPC Impact Story}
\textbf{Situation:}
WaveHPC supported multi-GPU distributed training workloads across shared cluster resources. As adoption increased, you observed inconsistent GPU utilization, long queue times, and inefficient scheduling for multi-node jobs.

\textbf{Task:}
Improve cluster efficiency and ensure fair, scalable execution for distributed training workloads across Slurm-based scheduling and Kubernetes-based services.

\textbf{Action:}
\begin{itemize}
\item Profiled cluster-wide GPU utilization with DGCM + Slurm metadata to identify idle GPUs during active allocations.
\item Analyzed job logs and NCCL traces to detect communication bottlenecks and straggler effects.
\item Identified workload imbalance across partitions and underutilized nodes.
\item Implemented automation pipelines for:
\begin{itemize}
\item job onboarding (correct \texttt{ntasks}, GPU binding, NCCL environment setup)
\item resource-aware Slurm submission defaults
\item monitoring dashboards for SM utilization, I/O wait, and network usage
\end{itemize}
\item Tuned job packing and recommended per-GPU batch scaling to reduce communication overhead.
\item Coordinated Slurm and Kubernetes CPU-pool segmentation for better multi-tenant isolation.
\end{itemize}

\textbf{Result:}
\begin{itemize}
\item improved cluster GPU utilization by \textbf{18\%}
\item reduced queue-time variance for multi-GPU jobs
\item decreased idle-GPU time during active allocations
\item established repeatable workload templates for distributed training
\end{itemize}

\textbf{Follow-up questions to expect:}
\begin{itemize}
\item What metric defined utilization? Primarily SM utilization, with memory utilization and network throughput for bottleneck context.
\item How did you separate compute underutilization from communication stalls? Low SM + high network implies communication bound; low SM + high CPU/I/O implies input bound.
\item How did you detect stragglers? Compared per-rank step times/utilization across ranks.
\item What happens when one rank is slower? Other GPUs block at synchronization; scaling efficiency drops.
\item How does small per-GPU batch hurt scaling? Less compute per step makes communication cost dominate.
\item How would you debug ~40\% GPU utilization? Check SM/memory utilization, NCCL time, DataLoader bottlenecks, network saturation, and stragglers.
\item Why were GPUs idle while allocated? Likely communication stalls, DataLoader bottlenecks, or incorrect GPU binding.
\item How prevent one tenant from monopolizing GPUs? Quotas + fairshare scheduling.
\end{itemize}

\subsection{High-Performance Matrix Inversion via Custom GPU Kernels}
\textbf{Situation:}
Matrix inversion on \(5\text{k}\times5\text{k}\) matrices was CPU-bound and became a bottleneck in numerical workloads, especially during repeated benchmarking.

\textbf{Task:}
Design a GPU-based implementation that improves runtime while preserving numerical stability and correctness, and understand what happens inside the "black box" beyond cuBLAS.

\textbf{Action:}
\begin{itemize}
\item Profiled the CPU baseline and identified heavy memory movement, serialization in Gaussian-elimination phases, and poor cache reuse.
\item Implemented block-based Gaussian elimination kernels in Triton and OpenCL for cross-vendor practice.
\item Applied tiling to increase arithmetic intensity.
\item Used shared/local memory to cache pivot rows and reduce global-memory pressure.
\item Carefully synchronized pivot updates across blocks.
\item Validated correctness with \(A\times A^{-1}=I\) checks and randomized vector tests.
\end{itemize}

\textbf{Result:}
\begin{itemize}
\item achieved \textbf{~1.5\(\times\)} speedup
\item reduced runtime by \textbf{<50\%} vs optimized CPU baseline
\item confirmed that memory layout and synchronization patterns dominated performance
\item gained practical intuition for memory hierarchy effects on throughput
\end{itemize}

\textbf{Follow-up questions to expect:}
\begin{itemize}
\item Was it compute-bound or memory-bound? Mostly memory-bound for large matrices.
\item How did you know? SM utilization was not maxed while memory throughput stayed high.
\item Why tiling? Better cache reuse and fewer global-memory loads.
\item What limited further speedup? Memory bandwidth and synchronization around pivot steps.
\item What is memory coalescing? Threads in a warp access contiguous memory so loads combine efficiently.
\item What is register pressure? Too many per-thread variables reduce occupancy and can cause spill to memory.
\end{itemize}

\subsection{GPU-Accelerated Image Processing Pipeline}
\textbf{Situation:}
High-resolution image filters (for example blur and sharpen) were CPU-bound, causing preprocessing bottlenecks.

\textbf{Task:}
Build a GPU-accelerated version with stable performance across different image sizes.

\textbf{Action:}
\begin{itemize}
\item Designed custom CUDA kernels in C++.
\item Tuned thread-block dimensions to maximize occupancy.
\item Used shared memory to cache image tiles.
\item Ensured coalesced pixel access to reduce global-memory latency.
\item Benchmarked against Python/PyTorch pipelines and validated correctness.
\item Compared kernel launch configs to minimize divergence.
\end{itemize}

\textbf{Result:}
\begin{itemize}
\item achieved up to \textbf{8\(\times\)} speedup vs CPU baseline
\item maintained pixel-level numerical parity
\item showed scaling with image resolution
\end{itemize}

\textbf{Follow-up questions to expect:}
\begin{itemize}
\item How did you choose block size? Empirical tuning (for example 16\(\times\)16 and 32\(\times\)32) to balance occupancy and shared-memory usage.
\item What is warp divergence? Threads in the same warp take different branches, reducing efficiency.
\item Was it memory-bound? Yes, pixel operations mainly stress memory bandwidth.
\item Why 8\(\times\) speedup? Massive parallelism + shared-memory caching + coalesced loads.
\end{itemize}

\subsection{Convolutional Layer Optimization with Triton}
\textbf{Situation:}
cuDNN hides convolution internals, making kernel-level bottlenecks hard to inspect directly.

\textbf{Task:}
Implement and optimize a 2D convolution kernel to deepen understanding of GPU memory hierarchy and execution mapping.

\textbf{Action:}
\begin{itemize}
\item Built custom 2D convolution kernels in Triton.
\item Experimented with tile sizes, block/thread mapping, shared-memory staging, and loop unrolling.
\item Compared runtime against PyTorch built-in convolution.
\item Analyzed memory-throughput vs arithmetic-intensity tradeoffs.
\end{itemize}

\textbf{Result:}
\begin{itemize}
\item achieved up to \textbf{~1.5\(\times\)} speedup on mid-size feature maps
\item developed intuition for memory-bound vs compute-bound regimes
\item understood how tiling affects cache reuse and where kernel fusion helps
\end{itemize}

\textbf{Follow-up questions to expect:}
\begin{itemize}
\item Why can convolution be memory-bound? Large feature maps drive high memory traffic relative to compute.
\item How does tiling help? Increases data reuse in shared memory and reduces global-memory loads.
\item What is kernel fusion? Combine multiple operations into one kernel to reduce memory traffic.
\end{itemize}

\end{document}
